\input{preamble}

% OK, commencer ici.
%
\begin{document}

\title{Champs algébriques}


\maketitle

\phantomsection
\label{section-phantom}

\tableofcontents

\section{Introduction}
\label{section-introduction}

\noindent
C'est ici que nous définissons les champs algébriques et formulons quelques
observations très élémentaires. La philosophie générale consistera à n'imposer
aucune condition de séparation et à ajouter les conditions nécessaires pour que
les lemmes, propositions et théorèmes soient vrais/démontrables. Ainsi, les notions
étudiées ici diffèrent légèrement de celles que l'on trouve ailleurs dans la littérature, par exemple
\cite{LM-B}.

\medskip\noindent
Ce chapitre ne constitue pas une introduction aux champs algébriques.
Pour une présentation informelle des champs algébriques, voir
Introduction aux champs algébriques, Section
\ref{stacks-introduction-section-introduction}.


\section{Conventions}
\label{section-conventions}

\noindent
Les conventions utilisées dans ce chapitre sont les mêmes que dans le
chapitre consacré aux espaces algébriques. Nous les rappelons ici par commodité.

\medskip\noindent
Nous travaillons dans un grand site fppf convenable $\Sch_{fppf}$
comme dans Topologies, Définition \ref{topologies-definition-big-fppf-site}.
Ainsi, sauf mention explicite du contraire, tous les schémas seront des objets
de $\Sch_{fppf}$. Nous examinons les changements qu'entraîne le remplacement du grand
site fppf dans la
Section \ref{section-change-big-site}.

\medskip\noindent
Nous travaillerons toujours relativement à une base $S$ contenue dans $\Sch_{fppf}$.
Nous travaillerons alors avec le grand site fppf $(\Sch/S)_{fppf}$, voir
Topologies, Définition \ref{topologies-definition-big-small-fppf}.
On retrouve le cas absolu en prenant
$S = \Spec(\mathbf{Z})$.

\medskip\noindent
Si $U, T$ sont des schémas sur $S$, nous notons
$U(T)$ l'ensemble de ses points à valeurs dans $T$ {\it au-dessus de} $S$.
En formule : $U(T) = \Mor_S(T, U)$.

\medskip\noindent
Notons que tout recouvrement fpqc est un épimorphisme effectif
universel, voir
Descente, Lemme \ref{descent-lemma-fpqc-universal-effective-epimorphisms}.
Par conséquent, la topologie sur $\Sch_{fppf}$
est moins fine que la topologie canonique et tous les préfaisceaux représentables
sont des faisceaux.








\section{Notations}
\label{section-notation}

\noindent
Nous employons les lettres $S, T, U, V, X, Y$ pour désigner des schémas.
Nous employons les lettres $\mathcal{X}, \mathcal{Y}, \mathcal{Z}$ pour désigner des
catégories (fibrées, fibrées en groupoïdes, champs, ...)
au-dessus de $(\Sch/S)_{fppf}$. Nous employons les lettres minuscules
$f$, $g$ pour les foncteurs tels que $f : \mathcal{X} \to \mathcal{Y}$
au-dessus de $(\Sch/S)_{fppf}$.
Nous employons les majuscules $F$, $G$, $H$ pour les espaces algébriques sur $S$, et plus
généralement pour les préfaisceaux d'ensembles sur $(\Sch/S)_{fppf}$.
(Dans les chapitres ultérieurs, nous emploierons de nouveau aussi $X$, $Y$, etc.
pour les espaces algébriques.)

\medskip\noindent
Ces choix visent à distinguer clairement les différents
types d'objets de ce chapitre, afin d'en établir les fondements.









\section{Catégories fibrées en groupoïdes représentables}
\label{section-representable}

\noindent
Soit $S$ un schéma contenu dans $\Sch_{fppf}$.
L'objet d'étude fondamental de ce chapitre sera une
catégorie fibrée en groupoïdes
$p : \mathcal{X} \to (\Sch/S)_{fppf}$, voir
Categories, Définition \ref{categories-definition-fibred-groupoids}.
Nous dirons souvent simplement « soit $\mathcal{X}$ une catégorie fibrée
en groupoïdes au-dessus de $(\Sch/S)_{fppf}$ » pour désigner
cette situation. Un $1$-morphisme $\mathcal{X} \to \mathcal{Y}$ de catégories
fibrées en groupoïdes au-dessus de $(\Sch/S)_{fppf}$ sera un $1$-morphisme
dans la $2$-catégorie des catégories fibrées en groupoïdes au-dessus de
$(\Sch/S)_{fppf}$, voir
Categories,
Définition \ref{categories-definition-categories-fibred-in-groupoids-over-C}.
Il s'agit simplement d'un foncteur $\mathcal{X} \to \mathcal{Y}$ au-dessus de
$(\Sch/S)_{fppf}$.
Rappelons qu'il s'agit en fait d'une $(2, 1)$-catégorie et que tous les $2$-produits fibrés
existent.

\medskip\noindent
Soit $\mathcal{X}$ une catégorie fibrée en groupoïdes au-dessus de
$(\Sch/S)_{fppf}$. Rappelons que $\mathcal{X}$
est dite {\it représentable} s'il existe un
schéma $U \in \Ob((\Sch/S)_{fppf})$ et une
équivalence
$$
j : \mathcal{X} \longrightarrow (\Sch/U)_{fppf}
$$
de catégories au-dessus de $(\Sch/S)_{fppf}$, voir
Categories,
Définition \ref{categories-definition-representable-fibred-category}.
Nous dirons parfois que $\mathcal{X}$ est
{\it représentable par un schéma} pour distinguer ce cas de celui où
$\mathcal{X}$ est représentable par un espace algébrique (voir
ci-dessous).

\medskip\noindent
Si $\mathcal{X}, \mathcal{Y}$ sont fibrées en groupoïdes et
représentables par $U, V$, alors on a
\begin{equation}
\label{equation-morphisms-schemes}
\Mor_{\textit{Cat}/(\Sch/S)_{fppf}}(\mathcal{X}, \mathcal{Y})
\Big/
2\text{-isomorphisme}
=
\Mor_{\Sch/S}(U, V)
\end{equation}
voir
Categories,
Lemme \ref{categories-lemma-morphisms-representable-fibred-categories}.
Plus précisément, tout $1$-morphisme $\mathcal{X} \to \mathcal{Y}$
donne lieu à un morphisme $U \to V$. Réciproquement, étant donné un morphisme
de schémas $U \to V$ au-dessus de $S$, il existe un $1$-morphisme
$\phi : \mathcal{X} \to \mathcal{Y}$ qui donne lieu à $U \to V$
et qui est unique à unique $2$-isomorphisme près.






\section{Le lemme de 2-Yoneda}
\label{section-2-yoneda}

\noindent
Soient $U \in \Ob((\Sch/S)_{fppf})$ et $\mathcal{X}$ une
catégorie fibrée en groupoïdes au-dessus de $(\Sch/S)_{fppf}$.
Nous utiliserons fréquemment le lemme de $2$-Yoneda, voir
Categories, Lemme \ref{categories-lemma-yoneda-2category}.
Techniquement, il affirme qu'il existe une équivalence de catégories
$$
\Mor_{\textit{Cat}/(\Sch/S)_{fppf}}(
(\Sch/U)_{fppf}, \mathcal{X})
\longrightarrow
\mathcal{X}_U, \quad
f \longmapsto f(U/U).
$$
Il affirme que les $1$-morphismes $(\Sch/U)_{fppf} \to \mathcal{X}$
correspondent aux objets $x$ de la catégorie fibre $\mathcal{X}_U$.
En effet, un $1$-morphisme $f : (\Sch/U)_{fppf} \to \mathcal{X}$
donne l'objet $x = f(U/U) \in \Ob(\mathcal{X}_U)$.
Réciproquement, étant donnés un choix d'images inverses pour $\mathcal{X}$ comme dans
Categories,
Définition \ref{categories-definition-pullback-functor-fibred-category},
et un objet $x$ de $\mathcal{X}_U$, on obtient un foncteur
$(\Sch/U)_{fppf} \to \mathcal{X}$ défini sur les objets par la règle
$$
(\varphi : V \to U) \longmapsto \varphi^*x
$$
Par abus de notation, nous employons
$x : (\Sch/U)_{fppf} \to \mathcal{X}$
pour désigner ce foncteur. Il vérifie bien $x(U/U) = x$
et, de plus, pour tout autre foncteur $f$ tel que $f(U/U) = x$, il existe
une unique $2$-isomorphisme $x \to f$. Autrement dit, le foncteur $x$
est déterminé par l'objet $x$ à unique $2$-isomorphisme près.

\medskip\noindent
Nous l'utiliserons désormais sans autre mention.





\section{Morphismes représentables de catégories fibrées en groupoïdes}
\label{section-representable-morphism}


\noindent
Soient $\mathcal{X}$, $\mathcal{Y}$ des catégories fibrées en groupoïdes
au-dessus de $(\Sch/S)_{fppf}$. Soit $f : \mathcal{X} \to \mathcal{Y}$
un {\it $1$-morphisme représentable}, voir
Categories, Définition
\ref{categories-definition-representable-map-categories-fibred-in-groupoids}.
Cela signifie que, pour tout $U \in \Ob((\Sch/S)_{fppf})$ et
tout $y \in \Ob(\mathcal{Y}_U)$, le $2$-produit fibré
$(\Sch/U)_{fppf} \times_{y, \mathcal{Y}} \mathcal{X}$
est représentable comme catégorie fibrée en groupoïdes au-dessus de
$(\Sch/U)_{fppf}$. Choisissons un objet représentant $f_y : V_y \to U$
et une équivalence
$$
(\Sch/V_y)_{fppf}
\longrightarrow
(\Sch/U)_{fppf} \times_{y, \mathcal{Y}} \mathcal{X}.
$$
Le morphisme $f_y$ correspond à la projection
$(\Sch/V_y)_{fppf} \to
(\Sch/U)_{fppf} \times_\mathcal{Y} \mathcal{Y}
\to (\Sch/U)_{fppf}$
Voir Section \ref{section-representable},
Équation (\ref{equation-morphisms-schemes}). Nous représentons ceci par le diagramme
\begin{equation}
\label{equation-representable}
\vcenter{
\xymatrix{
V_y \ar@{~>}[r] \ar[d]_{f_y} &
(\Sch/V_y)_{fppf} \ar[d] \ar[r] &
\mathcal{X} \ar[d]^f \\
U \ar@{~>}[r] &
(\Sch/U)_{fppf} \ar[r]^-y &
\mathcal{Y}
}
}
\end{equation}
où les flèches ondulées représentent le plongement de $2$-Yoneda.
Voici quelques lemmes sur cette notion, valables dans une grande généralité
(plus précisément, ils valent pour les catégories fibrées en groupoïdes au-dessus de toute
catégorie de base qui possède des produits fibrés).

\begin{lemma}
\label{lemma-morphism-schemes-gives-representable-transformation}
Soit $f : X \to Y$ un morphisme de $(\Sch/S)_{fppf}$.
Alors le $1$-morphisme induit par $f$
$$
(\Sch/X)_{fppf} \longrightarrow (\Sch/Y)_{fppf}
$$
est un $1$-morphisme représentable.
\end{lemma}

\begin{proof}
C'est formel et ne repose que sur le fait que
la catégorie $(\Sch/S)_{fppf}$ possède des produits fibrés.
\end{proof}

\begin{lemma}
\label{lemma-representable-morphism-equivalent}
Soit $S$ un objet de $\Sch_{fppf}$.
Considérons un diagramme $2$-commutatif
$$
\xymatrix{
\mathcal{X}' \ar[r] \ar[d]_{f'} & \mathcal{X} \ar[d]^f \\
\mathcal{Y}' \ar[r] & \mathcal{Y}
}
$$
de $1$-morphismes de catégories fibrées en groupoïdes au-dessus de
$(\Sch/S)_{fppf}$.
Supposons que les flèches horizontales soient des équivalences.
Alors $f$ est représentable si et seulement si $f'$ est représentable.
\end{lemma}

\begin{proof}
Omis.
\end{proof}

\begin{lemma}
\label{lemma-composition-representable-transformations}
Soit $S$ un schéma contenu dans $\Sch_{fppf}$.
Soient $\mathcal{X}, \mathcal{Y}, \mathcal{Z}$
des catégories fibrées en groupoïdes au-dessus de $(\Sch/S)_{fppf}$
Soient $f : \mathcal{X} \to \mathcal{Y}$, $g : \mathcal{Y} \to \mathcal{Z}$
des $1$-morphismes représentables. Alors
$$
g \circ f : \mathcal{X} \longrightarrow \mathcal{Z}
$$
est un $1$-morphisme représentable.
\end{lemma}

\begin{proof}
C'est entièrement formel et vaut dans toute catégorie.
\end{proof}

\begin{lemma}
\label{lemma-base-change-representable-transformations}
Soit $S$ un schéma contenu dans $\Sch_{fppf}$.
Soient $\mathcal{X}, \mathcal{Y}, \mathcal{Z}$
des catégories fibrées en groupoïdes au-dessus de $(\Sch/S)_{fppf}$
Soit $f : \mathcal{X} \to \mathcal{Y}$ un $1$-morphisme représentable.
Soit $g : \mathcal{Z} \to \mathcal{Y}$ un $1$-morphisme quelconque.
Considérons le diagramme de produit fibré
$$
\xymatrix{
\mathcal{Z} \times_{g, \mathcal{Y}, f} \mathcal{X} \ar[r]_-{g'} \ar[d]_{f'} &
\mathcal{X} \ar[d]^f \\
\mathcal{Z} \ar[r]^g & \mathcal{Y}
}
$$
Alors le changement de base $f'$ est un $1$-morphisme représentable.
\end{lemma}

\begin{proof}
C'est entièrement formel et vaut dans toute catégorie.
\end{proof}

\begin{lemma}
\label{lemma-product-representable-transformations}
Soit $S$ un schéma contenu dans $\Sch_{fppf}$.
Soient $\mathcal{X}_i, \mathcal{Y}_i$ des catégories fibrées en groupoïdes au-dessus de
$(\Sch/S)_{fppf}$, $i = 1, 2$.
Soient $f_i : \mathcal{X}_i \to \mathcal{Y}_i$, $i = 1, 2$
des $1$-morphismes représentables.
Alors
$$
f_1 \times f_2 :
\mathcal{X}_1 \times \mathcal{X}_2
\longrightarrow
\mathcal{Y}_1 \times \mathcal{Y}_2
$$
est un $1$-morphisme représentable.
\end{lemma}

\begin{proof}
Écrivons $f_1 \times f_2$ comme le composé
$\mathcal{X}_1 \times \mathcal{X}_2 \to
\mathcal{Y}_1 \times \mathcal{X}_2 \to
\mathcal{Y}_1 \times \mathcal{Y}_2$.
La première flèche est le changement de base de $f_1$ par l'application
$\mathcal{Y}_1 \times \mathcal{X}_2 \to \mathcal{Y}_1$, et la seconde flèche
est le changement de base de $f_2$ par l'application
$\mathcal{Y}_1 \times \mathcal{Y}_2 \to \mathcal{Y}_2$.
Ce lemme est donc une conséquence formelle des Lemmes
\ref{lemma-composition-representable-transformations}
et \ref{lemma-base-change-representable-transformations}.
\end{proof}



\section{Catégories fibrées en groupoïdes scindées}
\label{section-split}

\noindent
Soit $S$ un schéma contenu dans $\Sch_{fppf}$.
Rappelons qu'étant donné un « préfaisceau de groupoïdes »
$$
F : (\Sch/S)_{fppf}^{opp} \longrightarrow \textit{Groupoids}
$$
on obtient une catégorie fibrée en groupoïdes $\mathcal{S}_F$ au-dessus de
$(\Sch/S)_{fppf}$, voir
Categories, Exemple \ref{categories-example-functor-groupoids}.
Toute catégorie fibrée en groupoïdes isomorphe (!) à l'une d'elles
est appelée une {\it catégorie fibrée en groupoïdes scindée}.
Toute catégorie fibrée en groupoïdes est équivalente à une catégorie scindée.

\medskip\noindent
Si $F$ est un préfaisceau d'ensembles, alors $\mathcal{S}_F$ est
fibrée en ensembles, voir
Categories,
Définition \ref{categories-definition-category-fibred-sets},
et
Categories, Exemple \ref{categories-example-presheaf}.
La règle $F \mapsto \mathcal{S}_F$ est en un certain sens pleinement fidèle
sur les préfaisceaux, voir
Categories, Lemme \ref{categories-lemma-2-category-fibred-sets}.
Si $F, G$ sont des préfaisceaux, alors
$$
\mathcal{S}_{F \times G}
=
\mathcal{S}_F \times_{(\Sch/S)_{fppf}} \mathcal{S}_G
$$
et si $F \to H$ et $G \to H$ sont des applications de préfaisceaux d'ensembles, alors
$$
\mathcal{S}_{F \times_H G} =
\mathcal{S}_F \times_{\mathcal{S}_H} \mathcal{S}_G
$$
où les membres de droite sont des $2$-produits fibrés. Cela résulte immédiatement
des définitions, puisque les catégories fibres de
$\mathcal{S}_F, \mathcal{S}_G, \mathcal{S}_H$ ne possèdent que des morphismes identité.

\medskip\noindent
Un cas encore plus particulier est celui où $F = h_X$ est un préfaisceau
représentable. Dans ce cas, on a
$\mathcal{S}_{h_X} = (\Sch/X)_{fppf}$, voir
Categories,
Exemple \ref{categories-example-fibred-category-from-functor-of-points}.

\medskip\noindent
Nous emploierons désormais la notation $\mathcal{S}_F$ sans autre mention.





\section{Catégories fibrées en groupoïdes représentables par des espaces algébriques}
\label{section-representable-by-algebraic-spaces}

\noindent
La représentabilité par des espaces algébriques, que nous étudions dans cette section,
est une notion légèrement plus faible que la représentabilité.
On aurait pu éviter cette discussion si nous avions travaillé avec une catégorie
$\textit{Spaces}_{fppf}$ d'espaces algébriques plutôt qu'avec la catégorie
$\Sch_{fppf}$. Il nous semble toutefois naturel de considérer la
catégorie des schémas comme la collection naturelle d'« objets tests » au-dessus
desquels sont définies les catégories fibres d'un champ algébrique.

\medskip\noindent
Par analogie avec Categories, Définitions
\ref{categories-definition-representable-fibred-category},
nous posons la définition suivante.

\begin{definition}
\label{definition-representable-by-algebraic-space}
Soit $S$ un schéma contenu dans $\Sch_{fppf}$.
Une catégorie fibrée en groupoïdes $p : \mathcal{X} \to (\Sch/S)_{fppf}$
est dite {\it représentable par un espace algébrique sur $S$}
s'il existe un espace algébrique $F$ sur $S$ et une équivalence
$j : \mathcal{X} \to \mathcal{S}_F$
de catégories au-dessus de $(\Sch/S)_{fppf}$.
\end{definition}

\noindent
Nous poursuivons notre abus de notation en omettant l'équivalence $j$
chaque fois que nous rencontrons une telle situation.
Il résulte formellement de ce qui précède que, si $\mathcal{X}$ est
représentable (par un schéma), elle est représentable par un
espace algébrique. Voici l'analogue de
Categories,
Lemme \ref{categories-lemma-characterize-representable-fibred-category}.

\begin{lemma}
\label{lemma-characterize-representable-by-space}
Soit $S$ un schéma contenu dans $\Sch_{fppf}$.
Soit $p : \mathcal{X} \to (\Sch/S)_{fppf}$
une catégorie fibrée en groupoïdes.
Alors $\mathcal{X}$ est représentable par un espace algébrique sur $S$
si et seulement si les conditions suivantes sont satisfaites :
\begin{enumerate}
\item $\mathcal{X}$ est fibrée en setoïdes\footnote{Cela signifie qu'elle
est fibrée en groupoïdes et que les objets des catégories fibres
n'ont pas d'automorphismes non triviaux, voir Categories,
Définition \ref{categories-definition-category-fibred-setoids}.}, et
\item le préfaisceau $U \mapsto \Ob(\mathcal{X}_U)/\!\!\cong$ est
un espace algébrique.
\end{enumerate}
\end{lemma}

\begin{proof}
Omis, mais voir Categories,
Lemme \ref{categories-lemma-characterize-representable-fibred-category}.
\end{proof}

\noindent
Si $\mathcal{X}, \mathcal{Y}$ sont fibrées en groupoïdes et
représentables par des espaces algébriques $F, G$ sur $S$, alors on a
\begin{equation}
\label{equation-morphisms-spaces}
\Mor_{\textit{Cat}/(\Sch/S)_{fppf}}(\mathcal{X}, \mathcal{Y})
\Big/
2\text{-isomorphisme}
=
\Mor_{\Sch/S}(F, G)
\end{equation}
voir
Categories, Lemme \ref{categories-lemma-2-category-fibred-setoids}.
Plus précisément, tout $1$-morphisme $\mathcal{X} \to \mathcal{Y}$
donne lieu à un morphisme $F \to G$. Réciproquement, étant donné un morphisme
de faisceaux $F \to G$ au-dessus de $S$, il existe un $1$-morphisme
$\phi : \mathcal{X} \to \mathcal{Y}$ qui donne lieu à $F \to G$
et qui est unique à unique $2$-isomorphisme près.



\section{Morphismes représentables par des espaces algébriques}
\label{section-morphisms-representable-by-algebraic-spaces}

\noindent
Par analogie avec Categories, Définition
\ref{categories-definition-representable-map-categories-fibred-in-groupoids},
nous posons la définition suivante.

\begin{definition}
\label{definition-representable-by-algebraic-spaces}
Soit $S$ un schéma contenu dans $\Sch_{fppf}$.
Un $1$-morphisme $f : \mathcal{X} \to \mathcal{Y}$ de
catégories fibrées en groupoïdes au-dessus de $(\Sch/S)_{fppf}$
est dit {\it représentable par des espaces algébriques} si,
pour tout $U \in \Ob((\Sch/S)_{fppf})$
et tout $y : (\Sch/U)_{fppf} \to \mathcal{Y}$,
la catégorie fibrée en groupoïdes
$$
(\Sch/U)_{fppf} \times_{y, \mathcal{Y}} \mathcal{X}
$$
au-dessus de $(\Sch/U)_{fppf}$
est représentable par un espace algébrique sur $U$.
\end{definition}

\noindent
Choisissons un espace algébrique $F_y$ sur $U$ qui représente
$(\Sch/U)_{fppf} \times_{y, \mathcal{Y}} \mathcal{X}$.
On peut considérer $F_y$ comme un espace algébrique sur $S$
muni d'un morphisme canonique $f_y : F_y \to U$
au-dessus de $S$, voir
Spaces, Section \ref{spaces-section-change-base-scheme}.
Voici le diagramme
\begin{equation}
\label{equation-representable-by-algebraic-spaces}
\vcenter{
\xymatrix{
F_y \ar[d]_{f_y} &
(\Sch/U)_{fppf} \times_{y, \mathcal{Y}} \mathcal{X}
\ar@{~>}[l] \ar[d]_{\text{pr}_0} \ar[r]_-{\text{pr}_1} &
\mathcal{X} \ar[d]^f \\
U &
(\Sch/U)_{fppf} \ar@{~>}[l] \ar[r]^-y &
\mathcal{Y}
}
}
\end{equation}
où les flèches ondulées représentent la construction qui associe
à un champ fibré en setoïdes le faisceau associé des classes d'isomorphie
de ses objets. Le carré de droite est
$2$-commutatif et est un carré $2$-cartésien.

\medskip\noindent
Voici l'analogue de Categories,
Lemme \ref{categories-lemma-criterion-representable-map-stack-in-groupoids}.

\begin{lemma}
\label{lemma-criterion-map-representable-spaces-fibred-in-groupoids}
Soit $S$ un schéma contenu dans $\Sch_{fppf}$.
Soit $f : \mathcal{X} \to \mathcal{Y}$ un $1$-morphisme
de catégories fibrées en groupoïdes au-dessus de $(\Sch/S)_{fppf}$.
Pour que $f$ soit représentable par des espaces algébriques, il faut et il
suffit que les conditions suivantes soient satisfaites :
\begin{enumerate}
\item pour tout schéma $U/S$, le
foncteur $f_U : \mathcal{X}_U \longrightarrow \mathcal{Y}_U$
entre les catégories fibres est fidèle, et
\item pour tout $U$ et tout $y \in \Ob(\mathcal{Y}_U)$, le préfaisceau
$$
(h : V \to U)
\longmapsto
\{(x, \phi) \mid x \in \Ob(\mathcal{X}_V), \phi : h^*y \to f(x)\}/\cong
$$
est un espace algébrique sur $U$.
\end{enumerate}
Nous avons choisi ici des images inverses pour $\mathcal{Y}$.
\end{lemma}

\begin{proof}
Cela résulte de la description des catégories fibres des $2$-produits fibrés
$(\Sch/U)_{fppf} \times_{y, \mathcal{Y}} \mathcal{X}$ donnée dans
Categories, Lemme \ref{categories-lemma-identify-fibre-product},
combinée au
Lemme \ref{lemma-characterize-representable-by-space}.
\end{proof}

\noindent
Voici quelques lemmes sur cette notion, valables dans une grande généralité.

\begin{lemma}
\label{lemma-representable-by-spaces-morphism-equivalent}
Soit $S$ un objet de $\Sch_{fppf}$.
Considérons un diagramme $2$-commutatif
$$
\xymatrix{
\mathcal{X}' \ar[r] \ar[d]_{f'} & \mathcal{X} \ar[d]^f \\
\mathcal{Y}' \ar[r] & \mathcal{Y}
}
$$
de $1$-morphismes de catégories fibrées en groupoïdes au-dessus de
$(\Sch/S)_{fppf}$.
Supposons que les flèches horizontales soient des équivalences.
Alors $f$ est représentable par des espaces algébriques
si et seulement si $f'$ est représentable par des espaces algébriques.
\end{lemma}

\begin{proof}
Omis.
\end{proof}

\begin{lemma}
\label{lemma-morphism-spaces-gives-representable-by-spaces}
Soit $S$ un objet de $\Sch_{fppf}$.
Soit $f : \mathcal{X} \to \mathcal{Y}$
un $1$-morphisme de catégories fibrées en groupoïdes au-dessus de $S$.
Si $\mathcal{X}$ et $\mathcal{Y}$ sont représentables par des
espaces algébriques sur $S$, alors le $1$-morphisme $f$
est représentable par des espaces algébriques.
\end{lemma}

\begin{proof}
Omis. Cela ne repose que sur le fait que
la catégorie des espaces algébriques sur $S$ possède des produits fibrés,
voir Spaces, Lemme \ref{spaces-lemma-fibre-product-spaces}.
\end{proof}

\begin{lemma}
\label{lemma-map-presheaves-representable-by-algebraic-spaces}
Soit $S$ un objet de $\Sch_{fppf}$.
Soit $a : F \to G$ un morphisme de préfaisceaux d'ensembles sur $(\Sch/S)_{fppf}$.
Notons $a' : \mathcal{S}_F  \to \mathcal{S}_G$ le morphisme associé
de catégories fibrées en ensembles.
Alors $a$ est représentable par des espaces algébriques (voir
Bootstrap,
Définition \ref{bootstrap-definition-morphism-representable-by-spaces})
si et seulement si $a'$ est représentable par des espaces algébriques.
\end{lemma}

\begin{proof}
Omis.
\end{proof}

\begin{lemma}
\label{lemma-map-fibred-setoids-representable-algebraic-spaces}
Soit $S$ un objet de $\Sch_{fppf}$.
Soit $f : \mathcal{X} \to \mathcal{Y}$ un $1$-morphisme de
catégories fibrées en setoïdes au-dessus de $(\Sch/S)_{fppf}$.
Soient $F$, resp.\ $G$, le préfaisceau qui associe à $T$
l'ensemble des classes d'isomorphie d'objets de
$\mathcal{X}_T$, resp.\ $\mathcal{Y}_T$.
Soit $a : F \to G$ le morphisme de préfaisceaux correspondant à $f$.
Alors $a$ est représentable par des espaces algébriques (voir
Bootstrap,
Définition \ref{bootstrap-definition-morphism-representable-by-spaces})
si et seulement si $f$ est représentable par des espaces algébriques.
\end{lemma}

\begin{proof}
Omis. Indication : combiner les
Lemmes \ref{lemma-representable-by-spaces-morphism-equivalent}
et \ref{lemma-map-presheaves-representable-by-algebraic-spaces}.
\end{proof}

\begin{lemma}
\label{lemma-base-change-representable-by-spaces}
Soit $S$ un schéma contenu dans $\Sch_{fppf}$.
Soient $\mathcal{X}, \mathcal{Y}, \mathcal{Z}$
des catégories fibrées en groupoïdes au-dessus de $(\Sch/S)_{fppf}$.
Soit $f : \mathcal{X} \to \mathcal{Y}$ un $1$-morphisme
représentable par des espaces algébriques.
Soit $g : \mathcal{Z} \to \mathcal{Y}$ un $1$-morphisme quelconque.
Considérons le diagramme de produit fibré
$$
\xymatrix{
\mathcal{Z} \times_{g, \mathcal{Y}, f} \mathcal{X} \ar[r]_-{g'} \ar[d]_{f'} &
\mathcal{X} \ar[d]^f \\
\mathcal{Z} \ar[r]^g & \mathcal{Y}
}
$$
Alors le changement de base $f'$ est un $1$-morphisme représentable par des
espaces algébriques.
\end{lemma}

\begin{proof}
C'est formel.
\end{proof}

\begin{lemma}
\label{lemma-base-change-by-space-representable-by-space}
Soit $S$ un schéma contenu dans $\Sch_{fppf}$.
Soient $\mathcal{X}, \mathcal{Y}, \mathcal{Z}$
des catégories fibrées en groupoïdes au-dessus de $(\Sch/S)_{fppf}$
Soient $f : \mathcal{X} \to \mathcal{Y}$,
$g : \mathcal{Z} \to \mathcal{Y}$ des $1$-morphismes.
Supposons que
\begin{enumerate}
\item $f$ soit représentable par des espaces algébriques, et
\item $\mathcal{Z}$ soit représentable par un espace algébrique sur $S$.
\end{enumerate}
Alors le $2$-produit fibré
$\mathcal{Z} \times_{g, \mathcal{Y}, f} \mathcal{X}$
est représentable par un espace algébrique.
\end{lemma}

\begin{proof}
C'est une reformulation de
Bootstrap, Lemme \ref{bootstrap-lemma-representable-by-spaces-over-space}.
Notons d'abord que
$\mathcal{Z} \times_{g, \mathcal{Y}, f} \mathcal{X}$
est fibrée en setoïdes au-dessus de $(\Sch/S)_{fppf}$.
Elle est donc équivalente à $\mathcal{S}_F$ pour un certain préfaisceau
$F$ sur $(\Sch/S)_{fppf}$, voir
Categories, Lemme \ref{categories-lemma-setoid-fibres}.
De plus, soit $G$ un espace algébrique qui représente
$\mathcal{Z}$. Le $1$-morphisme
$\mathcal{Z} \times_{g, \mathcal{Y}, f} \mathcal{X} \to \mathcal{Z}$
est représentable par des espaces algébriques d'après le
Lemme \ref{lemma-base-change-representable-by-spaces}.
Et $\mathcal{Z} \times_{g, \mathcal{Y}, f} \mathcal{X} \to \mathcal{Z}$
correspond à un morphisme $F \to G$ d'après
Categories, Lemme \ref{categories-lemma-2-category-fibred-setoids}.
Alors $F \to G$ est représentable par des espaces algébriques d'après le
Lemme \ref{lemma-map-fibred-setoids-representable-algebraic-spaces}.
Par conséquent,
Bootstrap, Lemme \ref{bootstrap-lemma-representable-by-spaces-over-space}
implique que $F$ est un espace algébrique, comme souhaité.
\end{proof}

\noindent
Soient $S$, $\mathcal{X}$, $\mathcal{Y}$, $\mathcal{Z}$, $f$, $g$ comme dans le
Lemme \ref{lemma-base-change-by-space-representable-by-space}.
Soient $F$ et $G$ des espaces algébriques sur $S$ tels que
$F$ représente $\mathcal{Z} \times_{g, \mathcal{Y}, f} \mathcal{X}$
et $G$ représente $\mathcal{Z}$. Le $1$-morphisme
$f' : \mathcal{Z} \times_{g, \mathcal{Y}, f} \mathcal{X} \to \mathcal{Z}$
correspond à un morphisme $f' : F \to G$ d'espaces algébriques
d'après (\ref{equation-morphisms-spaces}).
On a donc le diagramme suivant
\begin{equation}
\label{equation-representable-by-algebraic-spaces-on-space}
\vcenter{
\xymatrix{
F \ar[d]_{f'} &
\mathcal{Z} \times_{g, \mathcal{Y}, f} \mathcal{X}
\ar@{~>}[l] \ar[d] \ar[r] &
\mathcal{X} \ar[d]^f \\
G &
\mathcal{Z} \ar@{~>}[l] \ar[r]^-g &
\mathcal{Y}
}
}
\end{equation}
où les flèches ondulées représentent la construction qui associe
à un champ fibré en setoïdes le faisceau associé des classes d'isomorphie
de ses objets.

\begin{lemma}
\label{lemma-composition-representable-by-spaces}
Soit $S$ un schéma contenu dans $\Sch_{fppf}$.
Soient $\mathcal{X}, \mathcal{Y}, \mathcal{Z}$
des catégories fibrées en groupoïdes au-dessus de $(\Sch/S)_{fppf}$.
Si $f : \mathcal{X} \to \mathcal{Y}$, $g : \mathcal{Y} \to \mathcal{Z}$
sont des $1$-morphismes représentables par des espaces algébriques, alors
$$
g \circ f : \mathcal{X} \longrightarrow \mathcal{Z}
$$
est un $1$-morphisme représentable par des espaces algébriques.
\end{lemma}

\begin{proof}
Cela résulte du
Lemme \ref{lemma-base-change-by-space-representable-by-space}.
Les détails sont omis.
\end{proof}

\begin{lemma}
\label{lemma-product-representable-by-spaces}
Soit $S$ un schéma contenu dans $\Sch_{fppf}$.
Soient $\mathcal{X}_i, \mathcal{Y}_i$ des catégories fibrées en groupoïdes au-dessus de
$(\Sch/S)_{fppf}$, $i = 1, 2$.
Soient $f_i : \mathcal{X}_i \to \mathcal{Y}_i$, $i = 1, 2$
des $1$-morphismes représentables par des espaces algébriques.
Alors
$$
f_1 \times f_2 :
\mathcal{X}_1 \times \mathcal{X}_2
\longrightarrow
\mathcal{Y}_1 \times \mathcal{Y}_2
$$
est un $1$-morphisme représentable par des espaces algébriques.
\end{lemma}

\begin{proof}
Écrivons $f_1 \times f_2$ comme le composé
$\mathcal{X}_1 \times \mathcal{X}_2 \to
\mathcal{Y}_1 \times \mathcal{X}_2 \to
\mathcal{Y}_1 \times \mathcal{Y}_2$.
La première flèche est le changement de base de $f_1$ par l'application
$\mathcal{Y}_1 \times \mathcal{X}_2 \to \mathcal{Y}_1$, et la seconde flèche
est le changement de base de $f_2$ par l'application
$\mathcal{Y}_1 \times \mathcal{Y}_2 \to \mathcal{Y}_2$.
Ce lemme est donc une conséquence formelle des Lemmes
\ref{lemma-composition-representable-by-spaces}
et \ref{lemma-base-change-representable-by-spaces}.
\end{proof}

\begin{lemma}
\label{lemma-get-a-stack}
\begin{reference}
Lemme figurant dans un courriel de Matthew Emerton daté du 15 juin 2016
\end{reference}
Soit $S$ un schéma contenu dans $\Sch_{fppf}$.
Soient $\mathcal{X} \to \mathcal{Z}$ et $\mathcal{Y} \to \mathcal{Z}$
des $1$-morphismes de catégories fibrées en groupoïdes au-dessus de $(\Sch/S)_{fppf}$.
Si $\mathcal{X} \to \mathcal{Z}$ est représentable par des espaces algébriques
et si $\mathcal{Y}$ est un champ en groupoïdes, alors
$\mathcal{X} \times_\mathcal{Z} \mathcal{Y}$ est un champ en groupoïdes.
\end{lemma}

\begin{proof}
La propriété, pour un morphisme, d'être représentable par des espaces algébriques
est préservée par changement de base
(Lemme \ref{lemma-base-change-by-space-representable-by-space}) ;
ainsi, en passant au changement de base
$\mathcal{X} \times_\mathcal{Z} \mathcal{Y}$ au-dessus de $\mathcal{Y}$,
on peut se ramener au cas d'un morphisme de catégories
fibrées en groupoïdes $\mathcal{X} \to \mathcal{Y}$
qui est représentable par des espaces algébriques et
dont le but est un champ en groupoïdes ; il s'agit alors de démontrer
que $\mathcal{X}$ est également un champ en groupoïdes.
Cela résulte de Stacks, Lemme
\ref{stacks-lemma-relative-sheaf-over-stack-is-stack},
dont les hypothèses sont satisfaites en vertu du
Lemme \ref{lemma-criterion-map-representable-spaces-fibred-in-groupoids}.
\end{proof}







\section{Propriétés des morphismes représentables par des espaces algébriques}
\label{section-representable-properties}

\noindent
Voici la définition qui permet de procéder.

\begin{definition}
\label{definition-relative-representable-property}
Soit $S$ un schéma appartenant à $\Sch_{fppf}$.
Soit $f : \mathcal{X} \to \mathcal{Y}$ un $1$-morphisme
de catégories fibrées en groupoïdes au-dessus de $(\Sch/S)_{fppf}$.
Supposons que $f$ soit représentable par des espaces algébriques.
Soit $\mathcal{P}$ une propriété des morphismes d'espaces algébriques qui
\begin{enumerate}
\item est préservée par tout changement de base, et
\item est locale pour la topologie fppf sur la base, voir
Descent on Spaces,
Définition \ref{spaces-descent-definition-property-morphisms-local}.
\end{enumerate}
Dans ce cas, on dit que $f$ a la {\it propriété $\mathcal{P}$} si, pour tout
$U \in \Ob((\Sch/S)_{fppf})$ et
tout $y \in \mathcal{Y}_U$, le morphisme d'espaces algébriques qui en résulte
$f_y : F_y \to U$, voir le
diagramme (\ref{equation-representable-by-algebraic-spaces}),
a la propriété $\mathcal{P}$.
\end{definition}

\noindent
Il importe de noter que nous n'utiliserons cette définition que pour des
propriétés de morphismes stables par changement de base et
locales pour la topologie fppf sur le but. Ce n'est
pas que la définition n'ait autrement aucun sens ; c'est plutôt
que nous pourrions vouloir donner une autre définition qui soit
mieux adaptée à la propriété envisagée.

\begin{lemma}
\label{lemma-property-morphism-equivalent}
Soit $S$ un objet de $\Sch_{fppf}$.
Soit $\mathcal{P}$ comme dans la
Définition \ref{definition-relative-representable-property}.
Considérons un diagramme $2$-commutatif
$$
\xymatrix{
\mathcal{X}' \ar[r] \ar[d]_{f'} & \mathcal{X} \ar[d]^f \\
\mathcal{Y}' \ar[r] & \mathcal{Y}
}
$$
de $1$-morphismes de catégories fibrées en groupoïdes au-dessus de
$(\Sch/S)_{fppf}$.
Supposons que les flèches horizontales soient des équivalences et que $f$
(ou, de manière équivalente, $f'$) soit représentable par des espaces algébriques.
Alors $f$ a $\mathcal{P}$ si et seulement si $f'$ a $\mathcal{P}$.
\end{lemma}

\begin{proof}
Cet énoncé a bien un sens en vertu du
Lemme \ref{lemma-representable-by-spaces-morphism-equivalent}.
Démonstration omise.
\end{proof}

\noindent
Voici un contrôle de cohérence.

\begin{lemma}
\label{lemma-map-presheaves-representable-by-spaces-transformation-property}
Soit $S$ un schéma appartenant à $\Sch_{fppf}$.
Soit $a : F \to G$ une application de préfaisceaux sur $(\Sch/S)_{fppf}$.
Soit $\mathcal{P}$ comme dans la
Définition \ref{definition-relative-representable-property}.
Supposons que $a$ soit représentable par des espaces algébriques.
Alors $a : F \to G$ a la propriété $\mathcal{P}$ (voir
Bootstrap, Définition \ref{bootstrap-definition-property-transformation})
si et seulement si le morphisme correspondant
$\mathcal{S}_F \to \mathcal{S}_G$ de catégories fibrées en groupoïdes
a la propriété $\mathcal{P}$.
\end{lemma}

\begin{proof}
L'énoncé du lemme a bien un sens en vertu du
Lemme \ref{lemma-map-presheaves-representable-by-algebraic-spaces}.
Démonstration omise.
\end{proof}

\begin{lemma}
\label{lemma-map-fibred-setoids-property}
Soit $S$ un objet de $\Sch_{fppf}$. Soit $\mathcal{P}$ comme dans la
Définition \ref{definition-relative-representable-property}.
Soit $f : \mathcal{X} \to \mathcal{Y}$ un $1$-morphisme de
catégories fibrées en setoïdes au-dessus de $(\Sch/S)_{fppf}$.
Soient $F$, resp.\ $G$, le préfaisceau qui associe à $T$
l'ensemble des classes d'isomorphie d'objets de
$\mathcal{X}_T$, resp.\ $\mathcal{Y}_T$.
Soit $a : F \to G$ l'application de préfaisceaux correspondant à $f$.
Alors $a$ a $\mathcal{P}$ si et seulement si $f$ a $\mathcal{P}$.
\end{lemma}

\begin{proof}
L'énoncé du lemme a bien un sens en vertu du
Lemme \ref{lemma-map-fibred-setoids-representable-algebraic-spaces}.
Le lemme résulte de la combinaison des
Lemmes \ref{lemma-property-morphism-equivalent}
et \ref{lemma-map-presheaves-representable-by-spaces-transformation-property}.
\end{proof}

\begin{lemma}
\label{lemma-composition-representable-transformations-property}
Soit $S$ un schéma appartenant à $\Sch_{fppf}$.
Soient $\mathcal{X}$, $\mathcal{Y}$, $\mathcal{Z}$ des catégories fibrées
en groupoïdes au-dessus de $(\Sch/S)_{fppf}$.
Soit $\mathcal{P}$ une propriété comme dans la
Définition \ref{definition-relative-representable-property}
qui est stable par composition.
Soient $f : \mathcal{X} \to \mathcal{Y}$ et
$g : \mathcal{Y} \to \mathcal{Z}$ des $1$-morphismes
représentables par des espaces algébriques.
Si $f$ et $g$ ont la propriété $\mathcal{P}$, il en est de même de
$g \circ f : \mathcal{X} \to \mathcal{Z}$.
\end{lemma}

\begin{proof}
L'énoncé du lemme a bien un sens en vertu du
Lemme \ref{lemma-composition-representable-by-spaces}.
Démonstration omise.
\end{proof}

\begin{lemma}
\label{lemma-base-change-representable-transformations-property}
Soit $S$ un schéma appartenant à $\Sch_{fppf}$.
Soient $\mathcal{X}, \mathcal{Y}, \mathcal{Z}$
des catégories fibrées en groupoïdes au-dessus de $(\Sch/S)_{fppf}$.
Soit $\mathcal{P}$ une propriété comme dans la
Définition \ref{definition-relative-representable-property}.
Soit $f : \mathcal{X} \to \mathcal{Y}$ un $1$-morphisme
représentable par des espaces algébriques.
Soit $g : \mathcal{Z} \to \mathcal{Y}$ un $1$-morphisme quelconque.
Considérons le diagramme de $2$-produit fibré
$$
\xymatrix{
\mathcal{Z} \times_{g, \mathcal{Y}, f} \mathcal{X} \ar[r]_-{g'} \ar[d]_{f'} &
\mathcal{X} \ar[d]^f \\
\mathcal{Z} \ar[r]^g & \mathcal{Y}
}
$$
Si $f$ a $\mathcal{P}$, alors le changement de base $f'$
a $\mathcal{P}$.
\end{lemma}

\begin{proof}
L'énoncé du lemme a bien un sens en vertu du
Lemme \ref{lemma-base-change-representable-by-spaces}.
Démonstration omise.
\end{proof}

\begin{lemma}
\label{lemma-descent-representable-transformations-property}
Soit $S$ un schéma appartenant à $\Sch_{fppf}$.
Soient $\mathcal{X}, \mathcal{Y}, \mathcal{Z}$
des catégories fibrées en groupoïdes au-dessus de $(\Sch/S)_{fppf}$.
Soit $\mathcal{P}$ une propriété comme dans la
Définition \ref{definition-relative-representable-property}.
Soit $f : \mathcal{X} \to \mathcal{Y}$ un $1$-morphisme
représentable par des espaces algébriques.
Soit $g : \mathcal{Z} \to \mathcal{Y}$ un $1$-morphisme quelconque.
Considérons le diagramme de produit fibré
$$
\xymatrix{
\mathcal{Z} \times_{g, \mathcal{Y}, f} \mathcal{X} \ar[r]_-{g'} \ar[d]_{f'} &
\mathcal{X} \ar[d]^f \\
\mathcal{Z} \ar[r]^g & \mathcal{Y}
}
$$
Supposons que, pour tout schéma $U$ et tout objet $x$ de $\mathcal{Y}_U$,
il existe un recouvrement fppf $\{U_i \to U\}$ tel que $x|_{U_i}$
appartienne à l'image essentielle du foncteur
$g : \mathcal{Z}_{U_i} \to \mathcal{Y}_{U_i}$.
Dans ce cas, si $f'$ a $\mathcal{P}$, alors $f$ a $\mathcal{P}$.
\end{lemma}

\begin{proof}
Démonstration omise. Indication : comparer avec la démonstration de
Spaces,
Lemme \ref{spaces-lemma-descent-representable-transformations-property}.
\end{proof}

\begin{lemma}
\label{lemma-product-representable-transformations-property}
Soit $S$ un schéma appartenant à $\Sch_{fppf}$.
Soit $\mathcal{P}$ une propriété comme dans la
Définition \ref{definition-relative-representable-property}
qui est stable par composition.
Soient $\mathcal{X}_i, \mathcal{Y}_i$ des catégories fibrées en groupoïdes au-dessus de
$(\Sch/S)_{fppf}$, $i = 1, 2$.
Soient $f_i : \mathcal{X}_i \to \mathcal{Y}_i$, $i = 1, 2$,
des $1$-morphismes représentables par des espaces algébriques.
Si $f_1$ et $f_2$ ont la propriété $\mathcal{P}$, il en est de même de
$
f_1 \times f_2 :
\mathcal{X}_1 \times \mathcal{X}_2
\to
\mathcal{Y}_1 \times \mathcal{Y}_2
$.
\end{lemma}

\begin{proof}
L'énoncé du lemme a bien un sens en vertu du
Lemme \ref{lemma-product-representable-by-spaces}.
Démonstration omise.
\end{proof}

\begin{lemma}
\label{lemma-representable-transformations-property-implication}
Soit $S$ un schéma appartenant à $\Sch_{fppf}$.
Soient $\mathcal{X}$, $\mathcal{Y}$ des catégories fibrées en groupoïdes
au-dessus de $(\Sch/S)_{fppf}$.
Soit $f : \mathcal{X} \to \mathcal{Y}$ un $1$-morphisme représentable
par des espaces algébriques.
Soient $\mathcal{P}$, $\mathcal{P}'$ des propriétés comme dans la
Définition \ref{definition-relative-representable-property}.
Supposons que, pour tout morphisme d'espaces algébriques $a : F \to G$,
on ait $\mathcal{P}(a) \Rightarrow \mathcal{P}'(a)$.
Si $f$ a la propriété $\mathcal{P}$, alors
$f$ a la propriété $\mathcal{P}'$.
\end{lemma}

\begin{proof}
C'est formel.
\end{proof}

\begin{lemma}
\label{lemma-open-fibred-category-is-full}
Soit $S$ un schéma appartenant à $\Sch_{fppf}$.
Soit $j : \mathcal X \to \mathcal Y$ un $1$-morphisme de
catégories fibrées en groupoïdes au-dessus de $(\Sch/S)_{fppf}$.
Supposons que $j$ soit représentable par des espaces algébriques et soit un monomorphisme
(voir la
Définition \ref{definition-relative-representable-property}
et
Descent on Spaces, Lemme
\ref{spaces-descent-lemma-descending-property-monomorphism}).
Alors $j$ est pleinement fidèle sur les catégories fibres.
\end{lemma}

\begin{proof}
Nous avons vu dans le
Lemme \ref{lemma-criterion-map-representable-spaces-fibred-in-groupoids}
que $j$ est fidèle sur les catégories fibres. Considérons un schéma $U$,
deux objets $u, v$ de $\mathcal{X}_U$ et un isomorphisme
$t : j(u) \to j(v)$ dans $\mathcal{Y}_U$. Il faut construire un
isomorphisme dans $\mathcal{X}_U$ entre $u$ et $v$.
Par le lemme de $2$-Yoneda (voir Section \ref{section-2-yoneda}),
nous considérons $u$, $v$ comme des $1$-morphismes
$u, v : (\Sch/U)_{fppf} \to \mathcal{X}$
et considérons le $2$-produit fibré
$$
(\Sch/U)_{fppf} \times_{j \circ v, \mathcal{Y}} \mathcal{X}.
$$
Par hypothèse, celui-ci est représentable par un espace algébrique
$F_{j \circ v}$ au-dessus de $U$, et le morphisme
$F_{j \circ v} \to U$ est un monomorphisme.
Mais puisque $(1_U, v, 1_{j(v)})$ donne un $1$-morphisme de
$(\Sch/U)_{fppf}$ dans le $2$-produit fibré affiché,
on voit que $F_{j \circ v} = U$ (on utilise ici
que si $V \to U$ est un monomorphisme d'espaces algébriques qui possède une
section, alors $V = U$). Par conséquent, le $1$-morphisme de projection sur
la première coordonnée
$$
(\Sch/U)_{fppf} \times_{j \circ v, \mathcal{Y}} \mathcal{X}
\to (\Sch/U)_{fppf}
$$
est une équivalence de catégories fibres.
Puisque $(1_U, u, t)$ et $(1_U, v, 1_{j(v)})$ donnent deux
objets de $((\Sch/U)_{fppf} \times_{j \circ v, \mathcal{Y}}
\mathcal{X})_U$ qui ont la même première coordonnée, il doit
exister entre eux un $2$-morphisme dans le $2$-produit fibré.
C'est, par définition, un morphisme $\tilde t : u \to v$ tel que
$j(\tilde t) = t$.
\end{proof}

\noindent
Voici une caractérisation des catégories fibrées en groupoïdes
dont la diagonale est représentable par des espaces algébriques.

\begin{lemma}
\label{lemma-representable-diagonal}
Soit $S$ un schéma appartenant à $\Sch_{fppf}$.
Soit $\mathcal{X}$ une catégorie fibrée en groupoïdes au-dessus de
$(\Sch/S)_{fppf}$. Les assertions suivantes sont équivalentes :
\begin{enumerate}
\item la diagonale $\mathcal{X} \to \mathcal{X} \times \mathcal{X}$
est représentable par des espaces algébriques,
\item pour tout schéma $U$ au-dessus de $S$ et tous
$x, y \in \Ob(\mathcal{X}_U)$, le faisceau
$\mathit{Isom}(x, y)$ est un espace algébrique au-dessus de $U$,
\item pour tout schéma $U$ au-dessus de $S$ et tout $x \in \Ob(\mathcal{X}_U)$,
le $1$-morphisme associé $x : (\Sch/U)_{fppf} \to \mathcal{X}$
est représentable par des espaces algébriques,
\item pour toute paire de schémas $T_1, T_2$ au-dessus de $S$ et tous
$x_i \in \Ob(\mathcal{X}_{T_i})$, $i = 1, 2$, le $2$-produit fibré
$(\Sch/T_1)_{fppf} \times_{x_1, \mathcal{X}, x_2}
(\Sch/T_2)_{fppf}$
est représentable par un espace algébrique,
\item pour toute catégorie fibrée en groupoïdes représentable $\mathcal{U}$
au-dessus de $(\Sch/S)_{fppf}$, tout $1$-morphisme
$\mathcal{U} \to \mathcal{X}$ est représentable par des espaces algébriques,
\item pour toute paire $\mathcal{T}_1, \mathcal{T}_2$ de catégories fibrées
en groupoïdes représentables au-dessus de $(\Sch/S)_{fppf}$ et tous
$1$-morphismes $x_i : \mathcal{T}_i \to \mathcal{X}$, $i = 1, 2$, le
$2$-produit fibré $\mathcal{T}_1 \times_{x_1, \mathcal{X}, x_2} \mathcal{T}_2$
est représentable par un espace algébrique,
\item pour toute catégorie fibrée en groupoïdes $\mathcal{U}$
au-dessus de $(\Sch/S)_{fppf}$ qui est
représentable par un espace algébrique, tout $1$-morphisme
$\mathcal{U} \to \mathcal{X}$ est représentable par des espaces algébriques,
\item pour toute paire $\mathcal{T}_1, \mathcal{T}_2$ de catégories fibrées
en groupoïdes au-dessus de $(\Sch/S)_{fppf}$ qui sont représentables
par des espaces algébriques, et tous $1$-morphismes
$x_i : \mathcal{T}_i \to \mathcal{X}$, le
$2$-produit fibré $\mathcal{T}_1 \times_{x_1, \mathcal{X}, x_2} \mathcal{T}_2$
est représentable par un espace algébrique.
\end{enumerate}
\end{lemma}

\begin{proof}
L'équivalence de (1) et (2) résulte du
Lemme \ref{stacks-lemma-isom-as-2-fibre-product} de Stacks
et des définitions.
Démontrons l'équivalence de (1) et (3).
Posons $\mathcal{C} = (\Sch/S)_{fppf}$ pour la catégorie de base.
Nous utiliserons certaines des observations de la démonstration de l'énoncé analogue
Categories, Lemme \ref{categories-lemma-representable-diagonal-groupoids}.
Nous emploierons le symbole $\cong$ au sens d'« équivalence de catégories fibrées
en groupoïdes au-dessus de $\mathcal{C} = (\Sch/S)_{fppf}$ ».
Supposons (1). Soient $U$ et $x$ comme dans (3). Pour tout schéma $V$
et tout $y \in \Ob(\mathcal{X}_V)$, on voit (comparer la référence ci-dessus) que
$$
\mathcal{C}/U
\times_{x, \mathcal{X}, y}
\mathcal{C}/V
\cong
(\mathcal{C}/U \times_S V)
\times_{(x, y), \mathcal{X} \times \mathcal{X}, \Delta}
\mathcal{X}
$$
qui est représentable par un espace algébrique par hypothèse. Réciproquement,
supposons (3). Considérons un schéma $U$ quelconque au-dessus de $S$ et une paire $(x, x')$
d'objets de $\mathcal{X}$ au-dessus de $U$. Il faut montrer que
$\mathcal{X} \times_{\Delta, \mathcal{X} \times \mathcal{X}, (x, x')} U$
est représentable par un espace algébrique. C'est clair, car
(comparer la référence ci-dessus)
$$
\mathcal{X}
\times_{\Delta, \mathcal{X} \times \mathcal{X}, (x, x')}
\mathcal{C}/U
\cong
(\mathcal{C}/U \times_{x, \mathcal{X}, x'} \mathcal{C}/U)
\times_{\mathcal{C}/U \times_S U, \Delta}
\mathcal{C}/U
$$
et le membre de droite est représentable par un espace algébrique par hypothèse
et parce que la catégorie des espaces algébriques au-dessus de $S$ admet des produits fibrés
et contient $U$ et $S$.

\medskip\noindent
Les équivalences
(3) $\Leftrightarrow$ (4),
(5) $\Leftrightarrow$ (6)
et
(7) $\Leftrightarrow$ (8)
sont formelles. Les équivalences
(3) $\Leftrightarrow$ (5) et
(4) $\Leftrightarrow$ (6)
résultent du
Lemme \ref{lemma-representable-by-spaces-morphism-equivalent}.
Supposons (3), et soit $\mathcal{U} \to \mathcal{X}$ comme dans (7).
Pour démontrer (7), il faut montrer que, pour tout schéma $V$ et tout $1$-morphisme
$y : (\Sch/V)_{fppf} \to \mathcal{X}$, le $2$-produit fibré
$(\Sch/V)_{fppf} \times_{y, \mathcal{X}} \mathcal{U}$
est représentable par un espace algébrique. La propriété (3) nous dit
que $y$ est représentable par des espaces algébriques ; par conséquent, le
Lemme \ref{lemma-base-change-by-space-representable-by-space}
donne l'assertion voulue. Enfin, (7) implique directement (3).
\end{proof}

\noindent
Dans la situation du lemme, pour tout $1$-morphisme
$x : (\Sch/U)_{fppf} \to \mathcal{X}$ comme dans le lemme, il est légitime
de dire que $x$ a la propriété $\mathcal{P}$, pour toute propriété comme dans la
Définition \ref{definition-relative-representable-property}.
Cela vaut en particulier pour
$\mathcal{P} = $ « surjectif »,
$\mathcal{P} = $ « lisse » et
$\mathcal{P} = $ « \'etale »,
voir
Descent on Spaces,
Lemmes \ref{spaces-descent-lemma-descending-property-surjective},
\ref{spaces-descent-lemma-descending-property-smooth} et
\ref{spaces-descent-lemma-descending-property-etale}.
Nous utiliserons ces trois cas dans les définitions
des champs algébriques ci-dessous.








\section{Champs en groupoïdes}
\label{section-stacks}

\noindent
Soit $S$ un schéma appartenant à $\Sch_{fppf}$.
Rappelons qu'une catégorie $p : \mathcal{X} \to (\Sch/S)_{fppf}$
au-dessus de $(\Sch/S)_{fppf}$ est appelée un
{\it champ en groupoïdes} (voir
Stacks, Définition \ref{stacks-definition-stack-in-groupoids})
si et seulement si
\begin{enumerate}
\item $p : \mathcal{X} \to (\Sch/S)_{fppf}$ est fibrée
en groupoïdes au-dessus de $(\Sch/S)_{fppf}$,
\item pour tout $U \in \Ob((\Sch/S)_{fppf})$ et
tous $x, y\in \Ob(\mathcal{X}_U)$, le préfaisceau
$\mathit{Isom}(x, y)$ est un faisceau sur le site $(\Sch/U)_{fppf}$, et
\item pour tout recouvrement $\mathcal{U} = \{U_i \to U\}$ de
$(\Sch/S)_{fppf}$, toutes les données de descente $(x_i, \phi_{ij})$
relatives à $\mathcal{U}$ sont effectives.
\end{enumerate}
Pour des exemples, voir
Examples of Stacks,
Section \ref{examples-stacks-section-examples-stacks-in-groupoids}
et suivantes.










\section{Champs algébriques}
\label{section-algebraic-stacks}

\noindent
Voici la définition d'un champ algébrique. Remarquons que la condition
(2) permet de donner un sens à la condition de la partie (3) selon laquelle
$(\Sch/U)_{fppf} \to \mathcal{X}$
est lisse et surjective ; voir la discussion qui suit le
Lemme \ref{lemma-representable-diagonal}.

\begin{definition}
\label{definition-algebraic-stack}
Soit $S$ un schéma de base appartenant à $\Sch_{fppf}$.
Un {\it champ algébrique au-dessus de $S$} est une catégorie
$$
p : \mathcal{X} \to (\Sch/S)_{fppf}
$$
au-dessus de $(\Sch/S)_{fppf}$ qui possède les propriétés suivantes :
\begin{enumerate}
\item La catégorie $\mathcal{X}$ est un champ en groupoïdes au-dessus de
$(\Sch/S)_{fppf}$.
\item La diagonale
$\Delta : \mathcal{X} \to \mathcal{X} \times \mathcal{X}$
est représentable par des espaces algébriques.
\item Il existe un schéma $U \in \Ob((\Sch/S)_{fppf})$
et un $1$-morphisme $(\Sch/U)_{fppf} \to \mathcal{X}$
qui est surjectif et lisse\footnote{Dans les chapitres ultérieurs, nous noterons
simplement ce morphisme $U \to \mathcal{X}$, comme il est d'usage dans la littérature. Une autre
bonne possibilité consisterait à formuler cette condition par l'existence d'une
catégorie fibrée en groupoïdes représentable $\mathcal{U}$ et d'un $1$-morphisme
lisse et surjectif $\mathcal{U} \to \mathcal{X}$.}.
\end{enumerate}
\end{definition}

\noindent
Il existe quelques différences avec d'autres définitions de la littérature.

\medskip\noindent
La première est que nous exigeons que $\mathcal{X}$ soit un champ en groupoïdes
pour la topologie fppf, tandis que de nombreuses références utilisent la topologie \'etale.
La topologie fppf nous semble en quelque sorte être la topologie naturelle
avec laquelle travailler. En définitive, la $2$-catégorie de champs algébriques
ainsi obtenue est la même. Ceci est expliqué dans
Criteria for Representability, Section \ref{criteria-section-stacks-etale}.

\medskip\noindent
La deuxième est que nous exigeons seulement que l'application diagonale de $\mathcal{X}$ soit
représentable par des espaces algébriques, alors que la plupart des références imposent d'autres
conditions. Notre point de vue consiste à essayer de démontrer un certain
nombre des résultats qui suivent en supposant seulement que la diagonale
de $\mathcal{X}$ soit représentable par des espaces algébriques, puis à ajouter simplement
une hypothèse supplémentaire chaque fois que cela est nécessaire. Cela présente en outre
l'avantage que tout espace algébrique (tel que défini dans
Spaces, Définition \ref{spaces-definition-algebraic-space})
donne lieu à un champ algébrique.

\medskip\noindent
La troisième est que certains articles exigent l'existence d'un
schéma $U$ et d'un morphisme surjectif et \'etale $U \to \mathcal{X}$.
Dans l'article fondateur \cite{DM}, où les champs algébriques furent introduits,
Deligne et Mumford ont utilisé cette définition et montré que
le champ de modules des courbes stables de genre $g > 1$ est un champ algébrique
qui admet un recouvrement \'etale par un schéma. Michael Artin, voir
\cite{ArtinVersal}, a compris que de nombreux
résultats naturels sur les champs algébriques s'étendent au cas où l'on
suppose seulement l'existence d'un recouvrement lisse par un schéma. D'où notre choix ci-dessus.
Pour distinguer les deux cas, la littérature emploie les termes « champ de Deligne-Mumford »
et « champ d'Artin ». Nous réserverons le terme
« champ d'Artin » à un usage ultérieur (insérer ici une référence future), et continuerons
d'employer « champ algébrique » ; en revanche, nous emploierons « champ de Deligne-Mumford »
pour désigner les champs algébriques qui admettent un recouvrement \'etale par un
schéma.

\begin{definition}
\label{definition-deligne-mumford}
Soit $S$ un schéma appartenant à $\Sch_{fppf}$.
Soit $\mathcal{X}$ un champ algébrique au-dessus de $S$.
On dit que $\mathcal{X}$ est un {\it champ de Deligne-Mumford} s'il existe
un schéma $U$ et un morphisme \'etale surjectif
$(\Sch/U)_{fppf} \to \mathcal{X}$.
\end{definition}

\noindent
Nous comparerons plus loin notre notion de champ de Deligne-Mumford avec
la notion définie dans l'article de Deligne et Mumford
(voir insérer ici une référence future).

\medskip\noindent
La catégorie des champs algébriques au-dessus de $S$ forme une $2$-catégorie.
En voici la définition précise.

\begin{definition}
\label{definition-morphism-algebraic-stacks}
Soit $S$ un schéma appartenant à $\Sch_{fppf}$.
La {\it $2$-catégorie des champs algébriques au-dessus de $S$} est la
sous-$2$-catégorie de la $2$-catégorie des catégories fibrées en
groupoïdes au-dessus de $(\Sch/S)_{fppf}$ (voir
Categories,
Définition \ref{categories-definition-categories-fibred-in-groupoids-over-C})
définie comme suit :
\begin{enumerate}
\item Ses objets sont les catégories fibrées en groupoïdes
au-dessus de $(\Sch/S)_{fppf}$ qui sont des champs algébriques au-dessus de $S$.
\item Ses $1$-morphismes $f : \mathcal{X} \to \mathcal{Y}$ sont
les foncteurs quelconques de catégories au-dessus de $(\Sch/S)_{fppf}$, comme dans
Categories, Définition \ref{categories-definition-categories-over-C}.
\item Ses $2$-morphismes sont les transformations entre foncteurs
au-dessus de $(\Sch/S)_{fppf}$, comme dans
Categories, Définition \ref{categories-definition-categories-over-C}.
\end{enumerate}
\end{definition}

\noindent
Autrement dit, cette $2$-catégorie est la sous-$2$-catégorie pleine de
$\textit{Cat}/(\Sch/S)_{fppf}$ dont les objets sont les champs algébriques.
Remarquons que tout $2$-morphisme est automatiquement un isomorphisme.
Il s'agit donc en fait d'une $(2, 1)$-catégorie et non pas seulement d'une $2$-catégorie.

\medskip\noindent
Nous verrons plus loin (insérer ici une référence future) que cette $2$-catégorie
admet des $2$-produits fibrés.

\medskip\noindent
Comme dans la remarque ci-dessus, la $2$-catégorie des champs algébriques au-dessus de $S$ est une
sous-$2$-catégorie pleine de la $2$-catégorie des catégories fibrées en groupoïdes
au-dessus de $(\Sch/S)_{fppf}$. Il se trouve qu'elle est stable par
équivalences. En voici l'énoncé précis.

\begin{lemma}
\label{lemma-equivalent}
Soit $S$ un schéma appartenant à $\Sch_{fppf}$.
Soient $\mathcal{X}$, $\mathcal{Y}$ des catégories au-dessus de $(\Sch/S)_{fppf}$.
Supposons que $\mathcal{X}$, $\mathcal{Y}$ soient équivalentes comme catégories au-dessus de
$(\Sch/S)_{fppf}$. Alors $\mathcal{X}$ est un champ algébrique si et
seulement si $\mathcal{Y}$ est un champ algébrique. De même, $\mathcal{X}$
est un champ de Deligne-Mumford si et seulement si $\mathcal{Y}$ est un champ de Deligne-Mumford.
\end{lemma}

\begin{proof}
Supposons que $\mathcal{X}$ soit un champ algébrique (resp.\ un champ de Deligne-Mumford). Le
Lemme \ref{stacks-lemma-stack-in-groupoids-equivalent} de Stacks
implique alors que $\mathcal{Y}$ est un champ en groupoïdes au-dessus de
$\Sch_{fppf}$. Choisissons une équivalence $f : \mathcal{X} \to \mathcal{Y}$
au-dessus de $\Sch_{fppf}$. Elle donne un diagramme $2$-commutatif
$$
\xymatrix{
\mathcal{X} \ar[r]_f \ar[d]_{\Delta_\mathcal{X}} &
\mathcal{Y} \ar[d]^{\Delta_\mathcal{Y}} \\
\mathcal{X} \times \mathcal{X} \ar[r]^{f \times f} &
\mathcal{Y} \times \mathcal{Y}
}
$$
dont les flèches horizontales sont des équivalences. Cela implique que
$\Delta_\mathcal{Y}$ est représentable par des espaces algébriques d'après le
Lemme \ref{lemma-representable-by-spaces-morphism-equivalent}.
Enfin, soit $U$ un schéma au-dessus de $S$, et soit
$x : (\Sch/U)_{fppf} \to \mathcal{X}$ un $1$-morphisme qui
est surjectif et lisse (resp.\ \'etale). En considérant le diagramme
$$
\xymatrix{
(\Sch/U)_{fppf} \ar[r]_{\text{id}} \ar[d]_x &
(\Sch/U)_{fppf} \ar[d]^{f \circ x} \\
\mathcal{X} \ar[r]^f &
\mathcal{Y}
}
$$
et en appliquant le
Lemme \ref{lemma-property-morphism-equivalent},
on conclut que $f \circ x$ est surjectif et lisse (resp.\ \'etale),
comme voulu.
\end{proof}





\section{Champs algébriques et espaces algébriques}
\label{section-stacks-spaces}

\noindent
Dans cette section, nous examinons quelques critères simples qui impliquent qu'un
champ algébrique est un espace algébrique. Le résultat principal affirme que c'est
exactement le cas lorsque les objets des catégories fibres n'ont aucun
automorphisme non trivial. Ce n'est pas une trivialité ! Avant d'y venir,
effectuons d'abord un contrôle de cohérence.

\begin{lemma}
\label{lemma-representable-algebraic}
Soit $S$ un schéma appartenant à $\Sch_{fppf}$.
\begin{enumerate}
\item Une catégorie fibrée en groupoïdes
$p : \mathcal{X} \to (\Sch/S)_{fppf}$
qui est représentable par un espace algébrique est un champ de Deligne-Mumford.
\item Si $F$ est un espace algébrique au-dessus de $S$, alors la
catégorie fibrée en groupoïdes associée
$p : \mathcal{S}_F \to (\Sch/S)_{fppf}$
est un champ de Deligne-Mumford.
\item Si $X \in \Ob((\Sch/S)_{fppf})$, alors
$(\Sch/X)_{fppf} \to (\Sch/S)_{fppf}$ est
un champ de Deligne-Mumford.
\end{enumerate}
\end{lemma}

\begin{proof}
Il est clair que (2) implique (3).
Les parties (1) et (2) sont équivalentes par le Lemme \ref{lemma-equivalent}.
Il suffit donc de démontrer (2).
Tout d'abord, remarquons que $\mathcal{S}_F$ est un champ en ensembles puisque
$F$ est un faisceau (Stacks, Lemme
\ref{stacks-lemma-stack-in-setoids-characterize}).
A fortiori, c'est un champ en groupoïdes. Ensuite, le morphisme diagonal
$\mathcal{S}_F \to \mathcal{S}_F  \times \mathcal{S}_F$
est le même que le morphisme $\mathcal{S}_F \to \mathcal{S}_{F \times F}$
qui provient de la diagonale de $F$. Il est donc représentable
par des espaces algébriques d'après le
Lemme \ref{lemma-morphism-spaces-gives-representable-by-spaces}.
En fait, il est même représentable (par des schémas), puisque la diagonale d'un
espace algébrique est représentable, mais nous n'en avons pas besoin.
Soit $U$ un schéma et soit $h_U \to F$ un morphisme \'etale surjectif.
Nous pouvons le considérer comme un morphisme \'etale surjectif d'espaces algébriques.
Ainsi, par le
Lemme
\ref{lemma-map-presheaves-representable-by-spaces-transformation-property},
le $1$-morphisme correspondant $(\Sch/U)_{fppf} \to \mathcal{S}_F$
est surjectif et \'etale.
\end{proof}

\noindent
Le résultat suivant affirme qu'un champ de Deligne-Mumford dont l'inertie
est triviale « est » un espace algébrique. Ce lemme sera rendu caduc par la
Proposition plus forte \ref{proposition-algebraic-stack-no-automorphisms}
ci-dessous, qui affirme que cela vaut plus généralement pour les champs algébriques…

\begin{lemma}
\label{lemma-algebraic-stack-no-automorphisms}
Soit $S$ un schéma appartenant à $\Sch_{fppf}$.
Soit $\mathcal{X}$ un champ algébrique au-dessus de $S$.
Les assertions suivantes sont équivalentes :
\begin{enumerate}
\item $\mathcal{X}$ est un champ de Deligne-Mumford et est un champ en setoïdes,
\item $\mathcal{X}$ est un champ de Deligne-Mumford tel que le
$1$-morphisme canonique $\mathcal{I}_\mathcal{X} \to \mathcal{X}$
soit une équivalence, et
\item $\mathcal{X}$ est représentable par un espace algébrique.
\end{enumerate}
\end{lemma}

\begin{proof}
L'équivalence de (1) et (2) résulte du
Lemme \ref{stacks-lemma-characterize-stack-in-setoids} de Stacks.
L'implication (3) $\Rightarrow$ (1) résulte du
Lemme \ref{lemma-representable-algebraic}.
Enfin, supposons (1). Par le
Lemme \ref{stacks-lemma-stack-in-setoids-characterize} de Stacks,
il existe un faisceau $F$ sur $(\Sch/S)_{fppf}$
et une équivalence $j : \mathcal{X} \to \mathcal{S}_F$. D'après le
Lemme \ref{lemma-map-presheaves-representable-by-algebraic-spaces},
le fait que $\Delta_\mathcal{X}$ soit représentable par des espaces
algébriques signifie que $\Delta_F : F \to F \times F$
est représentable par des espaces algébriques.
Soit $U$ un schéma et soit $x : (\Sch/U)_{fppf} \to \mathcal{X}$
un morphisme \'etale surjectif. La composée
$j \circ x : (\Sch/U)_{fppf} \to \mathcal{S}_F$
correspond à un morphisme $h_U \to F$ de faisceaux. Par le
Lemme \ref{bootstrap-lemma-representable-diagonal} de Bootstrap,
ce morphisme est représentable par des espaces algébriques.
Ainsi, par le
Lemme \ref{lemma-map-fibred-setoids-property},
on conclut que $h_U \to F$ est surjectif et \'etale.
Enfin, on applique le
Théorème \ref{bootstrap-theorem-bootstrap} de Bootstrap
pour voir que $F$ est un espace algébrique.
\end{proof}

\begin{proposition}
\label{proposition-algebraic-stack-no-automorphisms}
Soit $S$ un schéma appartenant à $\Sch_{fppf}$.
Soit $\mathcal{X}$ un champ algébrique au-dessus de $S$.
Les assertions suivantes sont équivalentes :
\begin{enumerate}
\item $\mathcal{X}$ est un champ en setoïdes,
\item le $1$-morphisme canonique $\mathcal{I}_\mathcal{X} \to \mathcal{X}$
est une équivalence, et
\item $\mathcal{X}$ est représentable par un espace algébrique.
\end{enumerate}
\end{proposition}

\begin{proof}
L'équivalence de (1) et (2) résulte du
Lemme \ref{stacks-lemma-characterize-stack-in-setoids} de Stacks.
L'implication (3) $\Rightarrow$ (1) résulte du
Lemme \ref{lemma-algebraic-stack-no-automorphisms}.
Enfin, supposons (1). Par le
Lemme \ref{stacks-lemma-stack-in-setoids-characterize} de Stacks,
il existe une équivalence $j : \mathcal{X} \to \mathcal{S}_F$
où $F$ est un faisceau sur $(\Sch/S)_{fppf}$.  D'après le
Lemme \ref{lemma-map-presheaves-representable-by-algebraic-spaces},
le fait que $\Delta_\mathcal{X}$ soit représentable par des espaces
algébriques signifie que $\Delta_F : F \to F \times F$
est représentable par des espaces algébriques.
Soit $U$ un schéma et soit $x : (\Sch/U)_{fppf} \to \mathcal{X}$
un morphisme lisse surjectif. La composée
$j \circ x : (\Sch/U)_{fppf} \to \mathcal{S}_F$
correspond à un morphisme $h_U \to F$ de faisceaux. Par le
Lemme \ref{bootstrap-lemma-representable-diagonal} de Bootstrap,
ce morphisme est représentable par des espaces algébriques.
Ainsi, par le
Lemme \ref{lemma-map-fibred-setoids-property},
on conclut que $h_U \to F$ est surjectif et lisse.
En particulier, il est surjectif, plat et localement de présentation finie
(par le
Lemme \ref{lemma-representable-transformations-property-implication}
et le fait qu'un morphisme lisse d'espaces algébriques est plat et
localement de présentation finie, voir
Morphisms of Spaces,
Lemmes \ref{spaces-morphisms-lemma-smooth-locally-finite-presentation} et
\ref{spaces-morphisms-lemma-smooth-flat}).
Enfin, on applique le
Théorème \ref{bootstrap-theorem-final-bootstrap} de Bootstrap
pour voir que $F$ est un espace algébrique.
\end{proof}







\section{2-produits fibrés de champs algébriques}
\label{section-2-fibre-products}

\noindent
La $2$-catégorie des champs algébriques admet des produits et des $2$-produits fibrés.
Le premier lemme est en réalité un cas particulier du
Lemme \ref{lemma-2-fibre-product},
mais sa démonstration est légèrement plus simple.

\begin{lemma}
\label{lemma-product-spaces}
Soit $S$ un schéma appartenant à $\Sch_{fppf}$.
Soient $\mathcal{X}$, $\mathcal{Y}$ des champs algébriques au-dessus de $S$.
Alors $\mathcal{X} \times_{(\Sch/S)_{fppf}} \mathcal{Y}$
est un champ algébrique, et est un produit dans la $2$-catégorie des
champs algébriques au-dessus de $S$.
\end{lemma}

\begin{proof}
Un objet de $\mathcal{X} \times_{(\Sch/S)_{fppf}} \mathcal{Y}$
au-dessus de $T$ est simplement une paire $(x, y)$, où $x$ est un objet de $\mathcal{X}_T$
et $y$ un objet de $\mathcal{Y}_T$. Il résulte donc immédiatement
des définitions que
$\mathcal{X} \times_{(\Sch/S)_{fppf}} \mathcal{Y}$ est un
champ en groupoïdes. Si $(x, y)$ et $(x', y')$ sont
deux objets de $\mathcal{X} \times_{(\Sch/S)_{fppf}} \mathcal{Y}$
au-dessus de $T$, alors
$$
\mathit{Isom}((x, y), (x', y')) =
\mathit{Isom}(x, x') \times \mathit{Isom}(y, y').
$$
Il résulte donc des équivalences du
Lemme \ref{lemma-representable-diagonal}
et du fait que la catégorie des espaces algébriques admet des produits
que la diagonale de $\mathcal{X} \times_{(\Sch/S)_{fppf}} \mathcal{Y}$
est représentable par des espaces algébriques.
Enfin, supposons que $U, V \in \Ob((\Sch/S)_{fppf})$,
et soient $x, y$ des morphismes lisses surjectifs
$x : (\Sch/U)_{fppf} \to \mathcal{X}$,
$y : (\Sch/V)_{fppf} \to \mathcal{Y}$.
Remarquons que
$$
(\Sch/U \times_S V)_{fppf} =
(\Sch/U)_{fppf}
\times_{(\Sch/S)_{fppf}} (\Sch/V)_{fppf}.
$$
L'objet $(\text{pr}_U^*x, \text{pr}_V^*y)$ de
$\mathcal{X} \times_{(\Sch/S)_{fppf}} \mathcal{Y}$ au-dessus de
$(\Sch/U \times_S V)_{fppf}$ définit ainsi un $1$-morphisme
$$
(\Sch/U \times_S V)_{fppf}
\longrightarrow
\mathcal{X} \times_{(\Sch/S)_{fppf}} \mathcal{Y}
$$
qui est la composée de changements de base de $x$ et de $y$, donc
est surjectif et lisse, voir les
Lemmes \ref{lemma-base-change-representable-transformations-property} et
\ref{lemma-composition-representable-transformations-property}.
On conclut que $\mathcal{X} \times_{(\Sch/S)_{fppf}} \mathcal{Y}$
est bien un champ algébrique. Nous omettons de vérifier qu'il
est effectivement un produit.
\end{proof}

\begin{lemma}
\label{lemma-2-fibre-product-general}
Soit $S$ un schéma appartenant à $\Sch_{fppf}$.
Soit $\mathcal{Z}$ un champ en groupoïdes au-dessus de $(\Sch/S)_{fppf}$
dont la diagonale est représentable par des espaces algébriques.
Soient $\mathcal{X}$, $\mathcal{Y}$ des champs algébriques au-dessus de $S$.
Soient $f : \mathcal{X} \to \mathcal{Z}$, $g : \mathcal{Y} \to \mathcal{Z}$
des $1$-morphismes de champs en groupoïdes. Alors le $2$-produit fibré
$\mathcal{X} \times_{f, \mathcal{Z}, g} \mathcal{Y}$ est un champ algébrique.
\end{lemma}

\begin{proof}
Il faut vérifier les conditions (1), (2) et (3) de la
Définition \ref{definition-algebraic-stack}.
La première condition résulte du
Lemme \ref{stacks-lemma-2-product-stacks-in-groupoids} de Stacks.

\medskip\noindent
La deuxième condition à vérifier est que les faisceaux $\mathit{Isom}$
soient représentables par des espaces algébriques. Pour cela, supposons que
$T$ soit un schéma au-dessus de $S$, et que $u, v$ soient des objets de
$(\mathcal{X} \times_{f, \mathcal{Z}, g} \mathcal{Y})_T$.
D'après notre construction des $2$-produits fibrés (qui remonte jusqu'à
Categories, Lemme \ref{categories-lemma-2-product-categories-over-C}),
nous pouvons écrire $u = (x, y, \alpha)$ et $v = (x', y', \alpha')$.
Ici, $\alpha : f(x) \to g(y)$, et de même pour $\alpha'$.
Il est alors clair que
$$
\xymatrix{
\mathit{Isom}(u, v) \ar[d] \ar[rr] & &
\mathit{Isom}(y, y') \ar[d]^{\phi \mapsto g(\phi) \circ \alpha} \\
\mathit{Isom}(x, x') \ar[rr]^-{\psi \mapsto \alpha' \circ f(\psi)} & &
\mathit{Isom}(f(x), g(y'))
}
$$
est un diagramme cartésien de faisceaux sur $(\Sch/T)_{fppf}$.
Puisque, par hypothèse, les faisceaux
$\mathit{Isom}(y, y')$, $\mathit{Isom}(x, x')$, $\mathit{Isom}(f(x), g(y'))$
sont des espaces algébriques (voir le
Lemme \ref{lemma-representable-diagonal}),
on voit que $\mathit{Isom}(u, v)$
est un espace algébrique.

\medskip\noindent
Soient $U, V \in \Ob((\Sch/S)_{fppf})$,
et soient $x, y$ des morphismes lisses surjectifs
$x : (\Sch/U)_{fppf} \to \mathcal{X}$,
$y : (\Sch/V)_{fppf} \to \mathcal{Y}$.
Considérons le morphisme
$$
(\Sch/U)_{fppf}
\times_{f \circ x, \mathcal{Z}, g \circ y}
(\Sch/V)_{fppf}
\longrightarrow
\mathcal{X} \times_{f, \mathcal{Z}, g} \mathcal{Y}.
$$
Comme la diagonale de $\mathcal{Z}$ est représentable par des espaces algébriques,
la source de cette flèche est représentable par un espace algébrique $F$, voir le
Lemme \ref{lemma-representable-diagonal}.
De plus, le morphisme est la composée
de changements de base de $x$ et de $y$, donc est surjectif et lisse, voir les
Lemmes \ref{lemma-base-change-representable-transformations-property} et
\ref{lemma-composition-representable-transformations-property}.
En choisissant un schéma $W$ et un morphisme \'etale surjectif $W \to F$,
on voit que la composée du $1$-morphisme affiché
avec le $1$-morphisme correspondant
$$
(\Sch/W)_{fppf}
\longrightarrow
(\Sch/U)_{fppf}
\times_{f \circ x, \mathcal{Z}, g \circ y}
(\Sch/V)_{fppf}
$$
est surjective et lisse, ce qui démontre la dernière condition.
\end{proof}

\begin{lemma}
\label{lemma-2-fibre-product}
Soit $S$ un schéma appartenant à $\Sch_{fppf}$.
Soient $\mathcal{X}, \mathcal{Y}, \mathcal{Z}$ des champs algébriques au-dessus de $S$.
Soient $f : \mathcal{X} \to \mathcal{Z}$, $g : \mathcal{Y} \to \mathcal{Z}$
des $1$-morphismes de champs algébriques. Alors le $2$-produit fibré
$\mathcal{X} \times_{f, \mathcal{Z}, g} \mathcal{Y}$ est un champ algébrique.
C'est aussi le $2$-produit fibré dans la $2$-catégorie des champs algébriques
au-dessus de $(\Sch/S)_{fppf}$.
\end{lemma}

\begin{proof}
Le fait que $\mathcal{X} \times_{f, \mathcal{Z}, g} \mathcal{Y}$ soit un
champ algébrique résulte du
Lemme plus fort \ref{lemma-2-fibre-product-general}.
Le fait que $\mathcal{X} \times_{f, \mathcal{Z}, g} \mathcal{Y}$
soit un $2$-produit fibré dans la $2$-catégorie des champs algébriques au-dessus de $S$
résulte formellement du fait que la $2$-catégorie des champs algébriques
au-dessus de $S$ est une sous-$2$-catégorie pleine de la $2$-catégorie des champs en
groupoïdes au-dessus de $(\Sch/S)_{fppf}$.
\end{proof}








\section{Champs algébriques, nouvelle présentation}
\label{section-overhaul}

\noindent
Quelques résultats fondamentaux sur les champs algébriques.

\begin{lemma}
\label{lemma-lift-morphism-presentations}
Soit $S$ un schéma appartenant à $\Sch_{fppf}$.
Soit $f : \mathcal{X} \to \mathcal{Y}$ un $1$-morphisme de champs
algébriques au-dessus de $S$.
Soit $V \in \Ob((\Sch/S)_{fppf})$.
Soit $y : (\Sch/V)_{fppf} \to \mathcal{Y}$ surjectif et lisse.
Alors il existe un objet $U \in \Ob((\Sch/S)_{fppf})$
et un diagramme $2$-commutatif
$$
\xymatrix{
(\Sch/U)_{fppf} \ar[r]_a \ar[d]_x &
(\Sch/V)_{fppf} \ar[d]^y \\
\mathcal{X} \ar[r]^f & \mathcal{Y}
}
$$
où $x$ est surjectif et lisse.
\end{lemma}

\begin{proof}
Choisissons d'abord $W \in \Ob((\Sch/S)_{fppf})$ et un $1$-morphisme
lisse surjectif $z : (\Sch/W)_{fppf} \to \mathcal{X}$.
Comme $\mathcal{Y}$ est un champ algébrique, on peut choisir une équivalence
$$
j :
\mathcal{S}_F
\longrightarrow
(\Sch/W)_{fppf}
\times_{f \circ z, \mathcal{Y}, y}
(\Sch/V)_{fppf}
$$
où $F$ est un espace algébrique. Par le
Lemme \ref{lemma-base-change-representable-transformations-property},
le morphisme
$\mathcal{S}_F \to (\Sch/W)_{fppf}$ est surjectif et lisse
comme changement de base de $y$. Ainsi, par le
Lemme \ref{lemma-composition-representable-transformations-property},
on voit que $\mathcal{S}_F \to \mathcal{X}$ est surjectif et lisse.
Choisissons un objet $U \in \Ob((\Sch/S)_{fppf})$
et un morphisme \'etale surjectif $U \to F$. En appliquant alors
une fois encore le Lemme \ref{lemma-composition-representable-transformations-property},
on obtient les propriétés voulues.
\end{proof}

\noindent
Ce lemme généralise la
Proposition \ref{proposition-algebraic-stack-no-automorphisms}.

\begin{lemma}
\label{lemma-characterize-representable-by-algebraic-spaces}
Soit $S$ un schéma appartenant à $\Sch_{fppf}$.
Soit $f : \mathcal{X} \to \mathcal{Y}$ un $1$-morphisme de champs
algébriques au-dessus de $S$. Les assertions suivantes sont équivalentes :
\begin{enumerate}
\item pour tout $U \in \Ob((\Sch/S)_{fppf})$,
le foncteur $f : \mathcal{X}_U \to \mathcal{Y}_U$ est fidèle,
\item le foncteur $f$ est fidèle, et
\item $f$ est représentable par des espaces algébriques.
\end{enumerate}
\end{lemma}

\begin{proof}
Les parties (1) et (2) sont équivalentes par les propriétés générales des $1$-morphismes
de catégories fibrées en groupoïdes, voir
Categories, Lemme \ref{categories-lemma-equivalence-fibred-categories}.
On voit que (3) implique (2) par le
Lemme \ref{lemma-criterion-map-representable-spaces-fibred-in-groupoids}.
Enfin, supposons (2).
Soit $U$ un schéma. Soit $y \in \Ob(\mathcal{Y}_U)$.
Il faut démontrer que
$$
\mathcal{W} = (\Sch/U)_{fppf} \times_{y, \mathcal{Y}} \mathcal{X}
$$
est représentable par un espace algébrique au-dessus de $U$. Puisque
$(\Sch/U)_{fppf}$ est un champ algébrique, il résulte du
Lemme \ref{lemma-2-fibre-product}
que $\mathcal{W}$ est un champ algébrique.
D'autre part, la description explicite des objets de $\mathcal{W}$
comme triplets $(V, x, \alpha : y(V) \to f(x))$ et le fait que $f$ soit
fidèle montrent que les catégories fibres de $\mathcal{W}$ sont des setoïdes. Par conséquent, la
Proposition \ref{proposition-algebraic-stack-no-automorphisms}
garantit que $\mathcal{W}$ est représentable par un espace algébrique.
\end{proof}

\begin{lemma}
\label{lemma-smooth-surjective-morphism-implies-algebraic}
Soit $S$ un schéma appartenant à $\Sch_{fppf}$.
Soit $u : \mathcal{U} \to \mathcal{X}$ un $1$-morphisme de
champs en groupoïdes au-dessus de $(\Sch/S)_{fppf}$. Si
\begin{enumerate}
\item $\mathcal{U}$ est représentable par un espace algébrique, et
\item $u$ est représentable par des espaces algébriques, surjectif et lisse,
\end{enumerate}
alors $\mathcal X$ est un champ algébrique au-dessus de $S$.
\end{lemma}

\begin{proof}
Il faut montrer que $\Delta : \mathcal{X} \to \mathcal{X} \times \mathcal{X}$
est représentable par des espaces algébriques, voir la
Définition \ref{definition-algebraic-stack}.
Étant donnés deux schémas $T_1$, $T_2$ au-dessus de $S$, notons
$\mathcal{T}_i = (\Sch/T_i)_{fppf}$ les catégories fibrées représentables
associées. Supposons donnés des $1$-morphismes
$f_i : \mathcal{T}_i \to \mathcal{X}$.
D'après le
Lemme \ref{lemma-representable-diagonal},
il suffit de démontrer que le $2$-produit
fibré $\mathcal{T}_1 \times_\mathcal{X} \mathcal{T}_2$
est représentable par un espace algébrique. Par le
Lemme de Stacks
\ref{stacks-lemma-2-fibre-product-stacks-in-setoids-over-stack-in-groupoids},
il s'agit en tout cas d'un champ en setoïdes. Ainsi,
$\mathcal{T}_1 \times_\mathcal{X} \mathcal{T}_2$ correspond
à un faisceau $F$ sur $(\Sch/S)_{fppf}$, voir
Stacks, Lemme \ref{stacks-lemma-stack-in-setoids-characterize}.
Soit $U$ l'espace algébrique qui représente $\mathcal{U}$.
Par hypothèse,
$$
\mathcal{T}_i' = \mathcal{U} \times_{u, \mathcal{X}, f_i} \mathcal{T}_i
$$
est représentable par un espace algébrique $T'_i$ au-dessus de $S$. Par conséquent,
$\mathcal{T}_1' \times_\mathcal{U} \mathcal{T}_2'$ est représentable
par l'espace algébrique $T'_1 \times_U T'_2$.
Considérons le diagramme commutatif
$$
\xymatrix{
&
\mathcal{T}_1 \times_{\mathcal X} \mathcal{T}_2 \ar[rr]\ar'[d][dd] & &
\mathcal{T}_1 \ar[dd] \\
\mathcal{T}_1' \times_\mathcal{U} \mathcal{T}_2' \ar[ur]\ar[rr]\ar[dd] & &
\mathcal{T}_1' \ar[ur]\ar[dd] \\
&
\mathcal{T}_2 \ar'[r][rr] & &
\mathcal X \\
\mathcal{T}_2' \ar[rr]\ar[ur] & &
\mathcal{U} \ar[ur] }
$$
Dans ce diagramme, le carré inférieur, le carré droit, le carré arrière et
le carré avant sont des $2$-produits fibrés. Un argument formel montre alors
que $\mathcal{T}_1' \times_\mathcal{U} \mathcal{T}_2' \to
\mathcal{T}_1 \times_{\mathcal X} \mathcal{T}_2$
est le « changement de base » de $\mathcal{U} \to \mathcal{X}$ ; plus précisément,
le diagramme
$$
\xymatrix{
\mathcal{T}_1' \times_\mathcal{U} \mathcal{T}_2' \ar[d] \ar[r] &
\mathcal{U} \ar[d] \\
\mathcal{T}_1 \times_{\mathcal X} \mathcal{T}_2 \ar[r] &
\mathcal{X}
}
$$
est un carré $2$-cartésien.
Ainsi, $T'_1 \times_U T'_2 \to F$ est représentable par des espaces algébriques,
lisse et surjectif, voir les
Lemmes \ref{lemma-map-fibred-setoids-representable-algebraic-spaces},
\ref{lemma-base-change-representable-by-spaces},
\ref{lemma-map-fibred-setoids-property} et
\ref{lemma-base-change-representable-transformations-property}.
Par conséquent, $F$ est un espace algébrique par le
Théorème \ref{bootstrap-theorem-final-bootstrap} de Bootstrap,
ce qui conclut.
\end{proof}

\noindent
Une application du
Lemme \ref{lemma-smooth-surjective-morphism-implies-algebraic}
est que ce qui est un espace algébrique au-dessus d'un champ algébrique
est un champ algébrique. C'est l'analogue du
Lemme \ref{bootstrap-lemma-representable-by-spaces-over-space} de Bootstrap.
En fait, il suffit de supposer que le morphisme
$\mathcal{X} \to \mathcal{Y}$ est « algébrique », comme nous le verrons dans
Criteria for Representability,
Lemme \ref{criteria-lemma-algebraic-morphism-to-algebraic}.

\begin{lemma}
\label{lemma-representable-morphism-to-algebraic}
Soit $S$ un schéma appartenant à $\Sch_{fppf}$.
Soit $\mathcal{X} \to \mathcal{Y}$ un morphisme de champs en groupoïdes
au-dessus de $(\Sch/S)_{fppf}$. Supposons que
\begin{enumerate}
\item $\mathcal{X} \to \mathcal{Y}$ soit représentable par des espaces algébriques, et
\item $\mathcal{Y}$ soit un champ algébrique au-dessus de $S$.
\end{enumerate}
Alors $\mathcal{X}$ est un champ algébrique au-dessus de $S$.
\end{lemma}

\begin{proof}
Soit $\mathcal{V} \to \mathcal{Y}$ un $1$-morphisme lisse surjectif
d'un champ en groupoïdes représentable vers $\mathcal{Y}$. Il existe par la
Définition \ref{definition-algebraic-stack}.
Alors le $2$-produit fibré
$\mathcal{U} = \mathcal{V} \times_{\mathcal Y} \mathcal X$
est représentable par un espace algébrique d'après le
Lemme \ref{lemma-base-change-by-space-representable-by-space}.
Le $1$-morphisme $\mathcal{U} \to \mathcal X$ est représentable par des espaces
algébriques, lisse et surjectif, voir les
Lemmes \ref{lemma-base-change-representable-by-spaces} et
\ref{lemma-base-change-representable-transformations-property}.
Par le
Lemme \ref{lemma-smooth-surjective-morphism-implies-algebraic},
on conclut que $\mathcal{X}$ est un champ algébrique.
\end{proof}

\begin{lemma}
\label{lemma-open-fibred-category-is-algebraic}
\begin{reference}
La suppression de l'hypothèse selon laquelle $j$ est un monomorphisme a été signalée
dans un courriel de Matthew Emerton daté du 15 juin 2016.
\end{reference}
Soit $S$ un schéma appartenant à $\Sch_{fppf}$.
Soit $j : \mathcal X \to \mathcal Y$ un $1$-morphisme de
catégories fibrées en groupoïdes au-dessus de $(\Sch/S)_{fppf}$.
Supposons que $j$ soit représentable par des espaces algébriques.
Alors, si $\mathcal{Y}$ est un champ en groupoïdes
(resp.\ un champ algébrique), il en est de même de $\mathcal{X}$.
\end{lemma}

\begin{proof}
L'assertion concernant les champs algébriques résulte de celle concernant les
champs en groupoïdes par le Lemme \ref{lemma-representable-morphism-to-algebraic}.
Si $j$ est représentable par des espaces algébriques, alors $j$ est
fidèle sur les catégories fibres et, pour chaque $U$ et chaque
$y \in \Ob(\mathcal{Y}_U)$, le préfaisceau
$$
(h : V \to U)
\longmapsto
\{(x, \phi) \mid x \in \Ob(\mathcal{X}_V), \phi : h^*y \to f(x)\}/\cong
$$
est un espace algébrique au-dessus de $U$. Voir le
Lemme \ref{lemma-criterion-map-representable-spaces-fibred-in-groupoids}.
En particulier, ce préfaisceau est un faisceau, et la conclusion résulte du
Lemme \ref{stacks-lemma-relative-sheaf-over-stack-is-stack} de Stacks.
\end{proof}



\section{D'un champ algébrique à une présentation}
\label{section-stack-to-presentation}

\noindent
Étant donné un champ algébrique au-dessus de $S$, on obtient un groupoïde en espaces algébriques
au-dessus de $S$ dont le champ quotient associé est le champ algébrique.

\medskip\noindent
Rappelons que si $(U, R, s, t, c)$ est un groupoïde en espaces algébriques au-dessus de $S$,
alors $[U/R]$ désigne le champ quotient associé à cette donnée, voir
Groupoids in Spaces,
Définition \ref{spaces-groupoids-definition-quotient-stack}.
En général, $[U/R]$ n'est {\bf pas} un champ algébrique. En particulier, le
champ $[U/R]$ qui apparaît dans le lemme suivant n'est en général pas
algébrique.

\begin{lemma}
\label{lemma-map-space-into-stack}
Soit $S$ un schéma appartenant à $\Sch_{fppf}$.
Soit $\mathcal{X}$ un champ algébrique au-dessus de $S$.
Soit $\mathcal{U}$ un champ algébrique au-dessus de $S$ qui
est représentable par un espace algébrique.
Soit $f : \mathcal{U} \to \mathcal{X}$ un 1-morphisme. Alors
\begin{enumerate}
\item le $2$-produit fibré
$\mathcal{R} = \mathcal{U} \times_{f, \mathcal{X}, f} \mathcal{U}$
est représentable par un espace algébrique,
\item il existe une équivalence canonique
$$
\mathcal{U} \times_{f, \mathcal{X}, f} \mathcal{U}
\times_{f, \mathcal{X}, f} \mathcal{U} =
\mathcal{R} \times_{\text{pr}_1, \mathcal{U}, \text{pr}_0} \mathcal{R},
$$
\item la projection $\text{pr}_{02}$ induit, via (2), un $1$-morphisme
$$
\text{pr}_{02} :
\mathcal{R} \times_{\text{pr}_1, \mathcal{U}, \text{pr}_0} \mathcal{R}
\longrightarrow
\mathcal{R}
$$
\item soient $U$, $R$ les espaces algébriques représentant
$\mathcal{U}, \mathcal{R}$, et soient $t, s : R \to U$ et
$c : R \times_{s, U, t} R \to R$ les morphismes correspondant
aux $1$-morphismes
$\text{pr}_0, \text{pr}_1 : \mathcal{R} \to \mathcal{U}$
et au
$\text{pr}_{02} :
\mathcal{R} \times_{\text{pr}_1, \mathcal{U}, \text{pr}_0} \mathcal{R} \to
\mathcal{R}$ ci-dessus ; alors le quintuplet $(U, R, s, t, c)$ est un groupoïde en
espaces algébriques au-dessus de $S$,
\item le morphisme $f$ induit un $1$-morphisme canonique
$f_{can} : [U/R] \to \mathcal{X}$
de champs en groupoïdes au-dessus de $(\Sch/S)_{fppf}$, et
\item le $1$-morphisme $f_{can} : [U/R] \to \mathcal{X}$ est pleinement fidèle.
\end{enumerate}
\end{lemma}

\begin{proof}
Démonstration de (1). Par définition, $\Delta_\mathcal{X}$ est représentable
par des espaces algébriques ; le
Lemme \ref{lemma-representable-diagonal}
s'applique donc et montre que $\mathcal{U} \to \mathcal{X}$ est représentable
par des espaces algébriques. Le résultat découle alors du
Lemme \ref{lemma-base-change-by-space-representable-by-space}.

\medskip\noindent
Soit $T$ un schéma au-dessus de $S$. Par construction du $2$-produit fibré (voir
Categories, Lemme \ref{categories-lemma-2-product-categories-over-C}),
on voit que les objets de la catégorie fibre $\mathcal{R}_T$
sont les triplets $(a, b, \alpha)$ où $a, b \in \Ob(\mathcal{U}_T)$
et où $\alpha : f(a) \to f(b)$
est un morphisme de la catégorie fibre $\mathcal{X}_T$.

\medskip\noindent
Démonstration de (2). L'équivalence s'obtient en appliquant successivement
Categories, Lemmes \ref{categories-lemma-associativity-2-fibre-product} et
\ref{categories-lemma-2-fibre-product-erase-factor}.
Identifions
$\mathcal{U} \times_\mathcal{X} \mathcal{U} \times_\mathcal{X} \mathcal{U}$
à
$(\mathcal{U} \times_\mathcal{X} \mathcal{U})
\times_\mathcal{X} \mathcal{U}$.
Si $T$ est un schéma au-dessus de $S$, alors, sur les catégories fibres au-dessus de $T$,
cette équivalence envoie l'objet
$((a, b, \alpha), c, \beta)$ du membre de gauche
sur l'objet $((a, b, \alpha), (b, c, \beta))$ du membre de droite.

\medskip\noindent
Démonstration de (3). Le $1$-morphisme $\text{pr}_{02}$ est construit dans la démonstration de
Categories, Lemme \ref{categories-lemma-triple-2-fibre-product-pr02}.
À l'aide de la description ci-dessus des objets de la catégorie fibre,
on voit que $((a, b, \alpha), (b, c, \beta))$
est envoyé sur $(a, c, \beta \circ \alpha)$.

\medskip\noindent
Malheureusement, cela n'est {\it pas compatible} avec nos conventions sur les
groupoïdes, selon lesquelles on a toujours $j = (t, s) : R \to U$, et l'on « considère »
un point à valeurs dans $T$, noté $r$, de $R$ comme un morphisme $r : s(r) \to t(r)$.
Cela n'affecte toutefois pas la démonstration de (4), puisque l'opposé d'un
groupoïde est un groupoïde. Mais, dans la démonstration de (5), cela explique
les inverses dans la formule affichée ci-dessous.

\medskip\noindent
Démonstration de (4). Rappelons que le faisceau $U$ est isomorphe au faisceau
$T \mapsto \Ob(\mathcal{U}_T)/\!\cong$, et
de même pour $R$, voir le
Lemme \ref{lemma-characterize-representable-by-space}.
Il résulte du
Lemme de Categories
\ref{categories-lemma-category-fibred-setoids-presheaves-products}
que cette description est compatible aux $2$-produits fibrés ;
on obtient donc une identification analogue de
$\mathcal{R} \times_{\text{pr}_1, \mathcal{U}, \text{pr}_0} \mathcal{R}$
et de $R \times_{s, U, t} R$.
Les morphismes $t, s : R \to U$ et $c : R \times_{s, U, t} R \to R$
s'obtiennent à partir de l'égalité générale (\ref{equation-morphisms-spaces}).
Explicitement, ces applications sont les transformations de foncteurs qui proviennent
de l'action de $\text{pr}_0$, $\text{pr}_0$, $\text{pr}_{02}$
sur les classes d'isomorphie d'objets des catégories fibres.
Ainsi, pour montrer que l'on obtient un groupoïde en espaces
algébriques, il suffit de montrer que, pour tout schéma $T$ au-dessus de $S$,
la structure
$$
(\Ob(\mathcal{U}_T)/\!\cong,
\Ob(\mathcal{R}_T)/\!\cong,
\text{pr}_1, \text{pr}_0, \text{pr}_{02})
$$
est un groupoïde, ce qui est clair d'après notre description ci-dessus des objets de
$\mathcal{R}_T$.

\medskip\noindent
Démonstration de (5). Nous appliquerons finalement
Groupoids in Spaces,
Lemme \ref{spaces-groupoids-lemma-quotient-stack-2-coequalizer}
pour obtenir le foncteur $[U/R] \to \mathcal{X}$.
Considérons le $1$-morphisme $f : \mathcal{U} \to \mathcal{X}$.
Nous avons une $2$-flèche $\tau : f \circ \text{pr}_1 \to f \circ \text{pr}_0$
par définition de $\mathcal{R}$ comme $2$-produit fibré.
En effet, sur un objet $(a, b, \alpha)$ de $\mathcal{R}$ au-dessus de $T$, c'est
l'application $\alpha^{-1} : b \to a$. Nous affirmons que
$$
\tau \circ \text{id}_{\text{pr}_{02}} =
(\tau \star \text{id}_{\text{pr}_0})
\circ
(\tau \star \text{id}_{\text{pr}_1}).
$$
Cette identité dit qu'étant donné un objet
$((a, b, \alpha), (b, c, \beta))$ de
$\mathcal{R} \times_{\text{pr}_1, \mathcal{U}, \text{pr}_0} \mathcal{R}$
au-dessus de $T$, la composée de
$$
\xymatrix{
c \ar[r]^{\beta^{-1}} & b \ar[r]^{\alpha^{-1}} & a
}
$$
est la même que la flèche $(\beta \circ \alpha)^{-1} : a \to c$. C'est
manifestement vrai, d'où l'affirmation. On voit ainsi que toutes les
hypothèses du
Lemme de Groupoids in Spaces
\ref{spaces-groupoids-lemma-quotient-stack-2-coequalizer}
sont satisfaites pour la structure
$(\mathcal{U}, \mathcal{R}, \text{pr}_0, \text{pr}_1, \text{pr}_{02})$
ainsi que le $1$-morphisme $f$ et le $2$-morphisme $\tau$.
Cependant, pour appliquer le lemme, il faut démontrer cela
pour la structure $(\mathcal{S}_U, \mathcal{S}_R, s, t, c)$
munie de morphismes appropriés.

\medskip\noindent
Il devrait exister un argument général de théorie abstraite des catégories
qui transfère ces données de l'une à l'autre, mais il semble
assez long. Nous employons plutôt l'astuce suivante.
Choisissons un quasi-inverse $j^{-1} : \mathcal{S}_U \to \mathcal{U}$
de l'équivalence canonique $j : \mathcal{U} \to \mathcal{S}_U$ qui provient
de $U(T) = \Ob(\mathcal{U}_T)/\!\!\cong$.
Cela signifie simplement que, pour tout schéma $T/S$ et tout
objet $a \in \mathcal{U}_T$, nous avons choisi un
élément particulier de sa classe d'isomorphie, à savoir $j^{-1}(j(a))$.
Grâce à $j^{-1}$, on peut donc considérer $\mathcal{S}_U$
comme une sous-catégorie de $\mathcal{U}$. Cette sous-catégorie une fois choisie,
on peut considérer les objets $(a, b, \alpha)$ de $\mathcal{R}_T$
tels que $a, b$ soient des objets de $(\mathcal{S}_U)_T$, c'est-à-dire tels
que $j^{-1}(j(a)) = a$ et $j^{-1}(j(b)) = b$. Il est alors clair qu'ils
forment une sous-catégorie de $\mathcal{R}$ qui s'envoie isomorphiquement
sur $\mathcal{S}_R$ par l'équivalence canonique
$\mathcal{R} \to \mathcal{S}_R$. De plus, ceci est manifestement compatible
avec la formation du $2$-produit fibré
$\mathcal{R} \times_{\text{pr}_1, \mathcal{U}, \text{pr}_0} \mathcal{R}$.
On voit donc que l'on peut simplement restreindre
$f$ à $\mathcal{S}_U$ et restreindre $\tau$ en une transformation
entre foncteurs $\mathcal{S}_R \to \mathcal{X}$. Il est dès lors clair que
l'égalité affichée du
Lemme de Groupoids in Spaces
\ref{spaces-groupoids-lemma-quotient-stack-2-coequalizer}
est satisfaite, puisqu'elle l'est même comme égalité de transformations de foncteurs
$\mathcal{R} \times_{\text{pr}_1, \mathcal{U}, \text{pr}_0} \mathcal{R}
\to \mathcal{X}$ avant restriction à la sous-catégorie
$\mathcal{S}_{R \times_{s, U, t} R}$.

\medskip\noindent
Ceci démontre que le
Lemme de Groupoids in Spaces
\ref{spaces-groupoids-lemma-quotient-stack-2-coequalizer}
s'applique et fournit le morphisme de champs recherché
$f_{can} : [U/R] \to \mathcal{X}$. Précisons brièvement comment
$f_{can}$ est défini dans ce cas particulier.
Sur un objet $a$ de $\mathcal{S}_U$ au-dessus de $T$,
on a $f_{can}(a) = f(a)$, où l'on considère
$\mathcal{S}_U \subset \mathcal{U}$ par le plongement choisi ci-dessus.
Si $a, b$ sont des objets de $\mathcal{S}_U$ au-dessus de $T$, alors un morphisme
$\varphi : a \to b$ dans $[U/R]$ est, par définition, un objet de la
forme $\varphi = (b, a, \alpha)$ de $\mathcal{R}$ au-dessus de $T$. (Noter
l'interversion.) Et la règle de la démonstration du
Lemme de Groupoids in Spaces
\ref{spaces-groupoids-lemma-quotient-stack-2-coequalizer}
est
\begin{equation}
\label{equation-on-morphisms}
f_{can}(\varphi) = \Big(f(a) \xrightarrow{\alpha^{-1}} f(b)\Big).
\end{equation}
Démonstration de (6). $[U/R]$ et $\mathcal{X}$ sont tous deux des champs.
Ainsi, étant donnés un schéma $T/S$ et des objets $a, b$ de $[U/R]$
au-dessus de $T$, on obtient une transformation de faisceaux fppf
$$
\mathit{Isom}(a, b) \longrightarrow \mathit{Isom}(f_{can}(a), f_{can}(b))
$$
sur $(\Sch/T)_{fppf}$. Il faut montrer qu'il s'agit d'un
isomorphisme. On peut travailler localement pour la topologie fppf sur $T$, et donc supposer que
$a, b$ proviennent de morphismes $a, b : T \to U$. Par le plongement
$\mathcal{S}_U \subset \mathcal{U}$ ci-dessus, on peut aussi considérer $a, b$ comme
des objets de $\mathcal{U}$ au-dessus de $T$. Dans
Groupoids in Spaces,
Lemme \ref{spaces-groupoids-lemma-quotient-stack-morphisms},
nous avons vu que le faisceau de gauche est représenté par l'espace algébrique
$$
R \times_{(t, s), U \times_S U, (b, a)} T
$$
au-dessus de $T$. D'autre part, le membre de droite est, par le
Lemme \ref{stacks-lemma-isom-as-2-fibre-product} de Stacks,
égal au faisceau associé au champ en setoïdes suivant :
$$
\mathcal{X}
\times_{\mathcal{X} \times \mathcal{X}, (f \circ b, f \circ a)} T =
\mathcal{X}
\times_{\mathcal{X} \times \mathcal{X}, (f, f)}
(\mathcal{U} \times \mathcal{U})
\times_{\mathcal{U} \times \mathcal{U}, (b, a)} T =
\mathcal{R}
\times_{(\text{pr}_0, \text{pr}_1), \mathcal{U} \times \mathcal{U}, (b, a)} T
$$
qui est représentable par le produit fibré affiché ci-dessus.
À ce stade, nous avons montré que les deux faisceaux $\mathit{Isom}$
sont isomorphes. Notre $1$-morphisme $f_{can} : [U/R] \to \mathcal{X}$ induit
cet isomorphisme sur les faisceaux $\mathit{Isom}$ d'après
l'Équation (\ref{equation-on-morphisms}).
\end{proof}

\noindent
On peut utiliser le lemme très abstrait précédent pour produire des
présentations.

\begin{lemma}
\label{lemma-stack-presentation}
Soit $S$ un schéma appartenant à $\Sch_{fppf}$.
Soit $\mathcal{X}$ un champ algébrique au-dessus de $S$.
Soit $U$ un espace algébrique au-dessus de $S$.
Soit $f : \mathcal{S}_U \to \mathcal{X}$ un morphisme lisse surjectif.
Soient $(U, R, s, t, c)$ le groupoïde en espaces algébriques
et $f_{can} : [U/R] \to \mathcal{X}$ le résultat de l'application du
Lemme \ref{lemma-map-space-into-stack}
à $U$ et à $f$. Alors
\begin{enumerate}
\item les morphismes $s$, $t$ sont lisses, et
\item le $1$-morphisme $f_{can} : [U/R] \to \mathcal{X}$
est une équivalence.
\end{enumerate}
\end{lemma}

\begin{proof}
Les morphismes $s, t$ sont lisses par les
Lemmes \ref{lemma-property-morphism-equivalent} et
\ref{lemma-map-presheaves-representable-by-spaces-transformation-property}.
Comme le $1$-morphisme $f$ est lisse et
surjectif, il est clair que, pour tout schéma $T$ et tout objet
$a \in \Ob(\mathcal{X}_T)$, il existe un morphisme lisse et surjectif
$T' \to T$ tel que $a|_T'$ provienne d'un objet de
$[U/R]_{T'}$. Puisque $f_{can} : [U/R] \to \mathcal{X}$
est pleinement fidèle, on en déduit que
$[U/R] \to \mathcal{X}$ est essentiellement surjectif, car
les données de descente sur les objets sont effectives des deux côtés, voir
Stacks, Lemme \ref{stacks-lemma-characterize-essentially-surjective-when-ff}.
\end{proof}

\begin{remark}
\label{remark-flat-fp-presentation}
Si l'on suppose seulement que le morphisme $f : \mathcal{S}_U \to \mathcal{X}$ du
Lemme \ref{lemma-stack-presentation}
est surjectif, plat et localement de présentation finie, alors
$f_{can} : [U/R] \to \mathcal{X}$ sera encore une
équivalence. Dans ce cas, les morphismes $s$, $t$ seront plats et
localement de présentation finie, mais ne seront bien sûr pas lisses en général.
\end{remark}

\noindent
Le Lemme \ref{lemma-stack-presentation}
suggère les définitions suivantes.

\begin{definition}
\label{definition-smooth-groupoid}
Soit $S$ un schéma. Soit $B$ un espace algébrique au-dessus de $S$.
Soit $(U, R, s, t, c)$ un groupoïde en espaces algébriques au-dessus de $B$.
On dit que $(U, R, s, t, c)$ est un {\it groupoïde lisse}\footnote{Cette terminologie
peut prêter quelque peu à confusion : elle n'implique pas que $[U/R]$ soit lisse
au-dessus de quoi que ce soit.}
si $s, t : R \to U$ sont des morphismes lisses d'espaces algébriques.
\end{definition}

\begin{definition}
\label{definition-presentation}
Soit $\mathcal{X}$ un champ algébrique au-dessus de $S$.
Une {\it présentation} de $\mathcal{X}$ est la donnée d'un groupoïde lisse
$(U, R, s, t, c)$ en espaces algébriques au-dessus de $S$, et d'une
équivalence $f : [U/R] \to \mathcal{X}$.
\end{definition}

\noindent
Nous avons vu ci-dessus que tout champ algébrique admet une présentation.
Notre tâche suivante consiste à montrer que tout groupoïde lisse en espaces
algébriques au-dessus de $S$ donne lieu à un champ algébrique.


\section{Le champ algébrique associé à un groupoïde lisse}
\label{section-smooth-groupoid-gives-algebraic-stack}

\noindent
Dans cette section, nous partons d'un groupoïde lisse en espaces algébriques
et montrons que le champ quotient associé est un champ algébrique.

\begin{lemma}
\label{lemma-diagonal-quotient-stack}
Soit $S$ un schéma appartenant à $\Sch_{fppf}$.
Soit $(U, R, s, t, c)$ un groupoïde en espaces algébriques au-dessus de $S$.
Alors la diagonale de $[U/R]$ est représentable par des espaces algébriques.
\end{lemma}

\begin{proof}
Il suffit de montrer que les faisceaux $\mathit{Isom}$ sont des espaces
algébriques, voir le
Lemme \ref{lemma-representable-diagonal}.
Cela résulte du
Lemme \ref{bootstrap-lemma-quotient-stack-isom} de Bootstrap.
\end{proof}

\begin{lemma}
\label{lemma-smooth-quotient-smooth-presentation}
Soit $S$ un schéma appartenant à $\Sch_{fppf}$.
Soit $(U, R, s, t, c)$ un groupoïde lisse en espaces algébriques au-dessus de $S$.
Alors le morphisme $\mathcal{S}_U \to [U/R]$ est lisse et surjectif.
\end{lemma}

\begin{proof}
Soit $T$ un schéma et soit $x : (\Sch/T)_{fppf} \to [U/R]$
un $1$-morphisme. Il faut montrer que la projection
$$
\mathcal{S}_U \times_{[U/R]} (\Sch/T)_{fppf}
\longrightarrow
(\Sch/T)_{fppf}
$$
est surjective et lisse. Nous savons déjà que le membre de gauche
est représentable par un espace algébrique $F$, voir les
Lemmes \ref{lemma-diagonal-quotient-stack} et
\ref{lemma-representable-diagonal}.
Il faut donc montrer que le morphisme correspondant $F \to T$
d'espaces algébriques est surjectif et lisse.
Puisqu'il s'agit de propriétés de morphismes d'espaces algébriques
qui sont locales pour la topologie fppf sur le but, nous
pouvons les vérifier localement pour la topologie fppf sur $T$. Par construction, il existe
un recouvrement fppf $\{T_i \to T\}$ de $T$ tel que
$x|_{(\Sch/T_i)_{fppf}}$ provienne d'un morphisme
$x_i : T_i \to U$. (Remarquons que $F \times_T T_i$ représente le
$2$-produit fibré $\mathcal{S}_U \times_{[U/R]} (\Sch/T_i)_{fppf}$,
de sorte que tout est compatible au changement de base par $T_i \to T$.)
On peut donc supposer que $x$ provient de $x : T \to U$.
Dans ce cas, on voit que
$$
\mathcal{S}_U \times_{[U/R]} (\Sch/T)_{fppf}
=
(\mathcal{S}_U \times_{[U/R]} \mathcal{S}_U)
\times_{\mathcal{S}_U} (\Sch/T)_{fppf}
=
\mathcal{S}_R \times_{\mathcal{S}_U} (\Sch/T)_{fppf}
$$
La première égalité résulte du
Lemme \ref{categories-lemma-2-fibre-product-erase-factor} de Categories,
et la seconde du
Lemme de Groupoids in Spaces
\ref{spaces-groupoids-lemma-quotient-stack-2-cartesian}.
Manifestement, le dernier $2$-produit fibré est représenté par l'espace
algébrique $F = R \times_{s, U, x} T$, et la projection
$R \times_{s, U, x} T \to T$ est lisse comme changement de base du
morphisme lisse d'espaces algébriques $s : R \to U$.
Elle est aussi surjective, car $s$ possède une section (à savoir l'identité
$e : U \to R$ du groupoïde).
Cela démontre le lemme.
\end{proof}

\noindent
Voici le résultat principal de cette section.

\begin{theorem}
\label{theorem-smooth-groupoid-gives-algebraic-stack}
Soit $S$ un schéma appartenant à $\Sch_{fppf}$.
Soit $(U, R, s, t, c)$ un groupoïde lisse en espaces algébriques au-dessus de $S$.
Alors le champ quotient $[U/R]$ est un champ algébrique au-dessus de $S$.
\end{theorem}

\begin{proof}
Vérifions les trois conditions de la
Définition \ref{definition-algebraic-stack}.
Par construction, $[U/R]$ est un champ en groupoïdes,
ce qui est la première condition.

\medskip\noindent
La deuxième condition résulte du résultat plus fort qu'est le
Lemme \ref{lemma-diagonal-quotient-stack}.

\medskip\noindent
Enfin, il faut montrer qu'il existe un schéma $W$ au-dessus de $S$
et un $1$-morphisme lisse surjectif
$(\Sch/W)_{fppf} \longrightarrow \mathcal{X}$.
Choisissons d'abord $W \in \Ob((\Sch/S)_{fppf})$ et un
morphisme \'etale surjectif $W \to U$. Remarquons que cela
donne un morphisme \'etale surjectif $\mathcal{S}_W \to \mathcal{S}_U$
de catégories fibrées en ensembles, voir le
Lemme
\ref{lemma-map-presheaves-representable-by-spaces-transformation-property}.
Bien entendu, $\mathcal{S}_W \to \mathcal{S}_U$ est alors aussi surjectif et
lisse, voir le
Lemme \ref{lemma-representable-transformations-property-implication}.
Ainsi, $\mathcal{S}_W \to \mathcal{S}_U \to [U/R]$ est surjectif
et lisse par combinaison des
Lemmes \ref{lemma-smooth-quotient-smooth-presentation} et
\ref{lemma-composition-representable-transformations-property}.
\end{proof}





\section{Changement de grand site}
\label{section-change-big-site}

\noindent
Dans cette section, nous examinons brièvement ce qui se passe lorsque l'on change de grand site.
En substance, on peut toujours agrandir à volonté le grand site ; on peut donc
supposer que tout ensemble de schémas que l'on souhaite considérer appartient au grand
site fppf au-dessus duquel on considère notre espace algébrique.
Nous invitons le lecteur à sauter cette section.

\medskip\noindent
Les images réciproques de champs sont définies dans
Stacks, Section \ref{stacks-section-inverse-image}.

\begin{lemma}
\label{lemma-change-big-site}
Supposons donnés de grands sites $\Sch_{fppf}$ et $\Sch'_{fppf}$.
Supposons que $\Sch_{fppf}$ soit contenu dans $\Sch'_{fppf}$,
voir Topologies, Section \ref{topologies-section-change-alpha}.
Soit $S$ un objet de $\Sch_{fppf}$.
Soit $f : (\Sch'/S)_{fppf} \to (\Sch/S)_{fppf}$ le morphisme
de sites correspondant au foncteur d'inclusion
$u : (\Sch/S)_{fppf} \to (\Sch'/S)_{fppf}$.
Soit $\mathcal{X}$ un champ en groupoïdes au-dessus de $(\Sch/S)_{fppf}$.
\begin{enumerate}
\item si $\mathcal{X}$ est représentable par un
$X \in \Ob((\Sch/S)_{fppf})$, alors
$f^{-1}\mathcal{X}$ est aussi représentable ; il est même représenté par le
même schéma $X$, considéré maintenant comme objet de $(\Sch'/S)_{fppf}$,
\item si $\mathcal{X}$ est représentable par
$F \in \Sh((\Sch/S)_{fppf})$, qui est
un espace algébrique, alors $f^{-1}\mathcal{X}$ est représentable
par l'espace algébrique $f^{-1}F$,
\item si $\mathcal{X}$ est un champ algébrique, alors $f^{-1}\mathcal{X}$
est un champ algébrique, et
\item si $\mathcal{X}$ est un champ de Deligne-Mumford, alors $f^{-1}\mathcal{X}$
est aussi un champ de Deligne-Mumford.
\end{enumerate}
\end{lemma}

\begin{proof}
Démontrons (3). Par le
Lemme \ref{lemma-stack-presentation},
on peut écrire $\mathcal{X} = [U/R]$ pour un groupoïde lisse
en espaces algébriques $(U, R, s, t, c)$. Par le
Lemme de Groupoids in Spaces
\ref{spaces-groupoids-lemma-quotient-stack-change-big-site},
on voit que $f^{-1}[U/R] = [f^{-1}U/f^{-1}R]$.
Bien entendu, $(f^{-1}U, f^{-1}R, f^{-1}s, f^{-1}t, f^{-1}c)$
est aussi un groupoïde lisse en espaces algébriques. Cela démontre (3).

\medskip\noindent
Les autres cas (1), (2), (4) signifient chacun que $\mathcal{X}$ possède
une présentation $[U/R]$ d'un type particulier, et se traduisent donc par le
même type de présentation de $f^{-1}\mathcal{X} = [f^{-1}U/f^{-1}R]$.
Le lemme est ainsi démontré.
\end{proof}

\noindent
Il n'est pas vrai, en général, que la restriction d'un espace algébrique
sur le grand site le plus gros soit un espace algébrique sur le plus petit (pour de simples
raisons de cardinalité). On ne peut donc utiliser un lemme simple de ce
genre que pour agrandir la catégorie de base, jamais pour la rétrécir.

\begin{lemma}
\label{lemma-fully-faithful}
Supposons que $\Sch_{fppf}$ soit contenu dans $\Sch'_{fppf}$.
Soit $S$ un objet de $\Sch_{fppf}$. Notons
$\textit{Algebraic-Stacks}/S$ la $2$-catégorie des champs algébriques au-dessus de $S$
définie à l'aide de $\Sch_{fppf}$. De même, notons
$\textit{Algebraic-Stacks}'/S$ la $2$-catégorie des champs algébriques au-dessus de $S$
définie à l'aide de $\Sch'_{fppf}$. La règle
$\mathcal{X} \mapsto f^{-1}\mathcal{X}$ du
Lemme \ref{lemma-change-big-site}
définit un foncteur de $2$-catégories
$$
\textit{Algebraic-Stacks}/S \longrightarrow \textit{Algebraic-Stacks}'/S
$$
qui définit des équivalences de catégories de morphismes
$$
\Mor_{\textit{Algebraic-Stacks}/S}(\mathcal{X}, \mathcal{Y})
\longrightarrow
\Mor_{\textit{Algebraic-Stacks}'/S}(f^{-1}\mathcal{X}, f^{-1}\mathcal{Y})
$$
pour tous objets $\mathcal{X}, \mathcal{Y}$ de
$\textit{Algebraic-Stacks}/S$. Un objet
$\mathcal{X}'$ de $\textit{Algebraic-Stacks}'/S$
est équivalent à $f^{-1}\mathcal{X}$ pour un certain
$\mathcal{X}$ de $\textit{Algebraic-Stacks}/S$
si et seulement s'il possède une présentation $\mathcal{X} = [U'/R']$
où $U', R'$ sont isomorphes à $f^{-1}U$, $f^{-1}R$ pour certains
$U, R \in \textit{Spaces}/S$.
\end{lemma}

\begin{proof}
L'assertion concernant les catégories de morphismes résulte de l'énoncé plus général
Stacks, Lemme \ref{stacks-lemma-bigger-site}.
La caractérisation de l'« image essentielle » découle de la description
de $f^{-1}$ dans la démonstration du
Lemme \ref{lemma-change-big-site}.
\end{proof}


\section{Changement de schéma de base}
\label{section-change-base-scheme}

\noindent
Dans cette section, nous examinons brièvement ce qui se passe lorsque l'on change de schéma de base.
En substance, étant donné un morphisme $S \to S'$ de schémas de base, tout champ
algébrique au-dessus de $S$ peut être considéré comme un champ algébrique au-dessus de $S'$.

\begin{lemma}
\label{lemma-category-of-spaces-over-smaller-base-scheme}
Soit $\Sch_{fppf}$ un grand site fppf.
Soit $S \to S'$ un morphisme de ce site.
Les constructions A et B de
Stacks, Section \ref{stacks-section-localize}
ci-dessus donnent des isomorphismes de $2$-catégories
$$
\left\{
\begin{matrix}
2\text{-catégorie des champs}\\
\text{algébriques }\mathcal{X}\text{ au-dessus de }S
\end{matrix}
\right\}
\leftrightarrow
\left\{
\begin{matrix}
2\text{-catégorie des paires }(\mathcal{X}', f)\text{ formées d'un}\\
\text{champ algébrique }\mathcal{X}'\text{ au-dessus de }S'\text{ et d'un morphisme}\\
f : \mathcal{X}' \to (\Sch/S)_{fppf}\text{ de champs algébriques au-dessus de }S'
\end{matrix}
\right\}
$$
\end{lemma}

\begin{proof}
L'énoncé a bien un sens, car le foncteur
$j : (\Sch/S)_{fppf} \to (\Sch/S')_{fppf}$
est le foncteur de localisation associé à l'objet $S/S'$
de $(\Sch/S')_{fppf}$. Par le
Lemme \ref{stacks-lemma-localize-stacks} de Stacks,
il suffit de montrer que les constructions A et B
préservent les sous-catégories de champs algébriques.
Par exemple, si $\mathcal{X} = [U/R]$, alors la construction A
appliquée à $\mathcal{X}$ produit simplement
$\mathcal{X}' = \mathcal{X}$. Réciproquement, si $\mathcal{X}' = [U'/R']$,
le morphisme $p$ induit des morphismes d'espaces algébriques
$U' \to S$ et $R' \to S$, et alors $\mathcal{X} = [U'/R']$,
mais considéré maintenant comme un champ au-dessus de $S$. Le lemme est donc clair.
\end{proof}

\begin{definition}
\label{definition-viewed-as}
Soit $\Sch_{fppf}$ un grand site fppf.
Soit $S \to S'$ un morphisme de ce site.
Si $p : \mathcal{X} \to (\Sch/S)_{fppf}$
est un champ algébrique au-dessus de $S$, alors
$\mathcal{X}$ {\it considéré comme champ algébrique au-dessus de $S'$}
est le champ algébrique
$$
\mathcal{X} \longrightarrow (\Sch/S')_{fppf}
$$
obtenu en appliquant la construction A du
Lemme \ref{lemma-category-of-spaces-over-smaller-base-scheme}
à $\mathcal{X}$.
\end{definition}

\noindent
Réciproquement, que se passe-t-il si l'on part d'un champ algébrique $\mathcal{X}'$
au-dessus de $S'$ et que l'on souhaite obtenir un champ algébrique au-dessus de $S$ ?
On considère alors le $2$-produit fibré
$$
\mathcal{X}'_S
=
(\Sch/S)_{fppf} \times_{(\Sch/S')_{fppf}} \mathcal{X}'
$$
qui est un champ algébrique au-dessus de $S'$ d'après le
Lemme \ref{lemma-2-fibre-product}.
De plus, il est muni d'un $1$-morphisme naturel
$p : \mathcal{X}'_S \to (\Sch/S)_{fppf}$ et, par le
Lemme \ref{lemma-category-of-spaces-over-smaller-base-scheme},
il correspond donc canoniquement à un champ algébrique au-dessus de $S$.

\begin{definition}
\label{definition-change-of-base}
Soit $\Sch_{fppf}$ un grand site fppf.
Soit $S \to S'$ un morphisme de ce site.
Soit $\mathcal{X}'$ un champ algébrique au-dessus de $S'$.
Le {\it changement de base de $\mathcal{X}'$} est le
champ algébrique $\mathcal{X}'_S$ au-dessus de $S$ décrit ci-dessus.
\end{definition}




\begin{multicols}{2}[\section{Autres chapitres}]
\noindent
Préliminaires
\begin{enumerate}
\item \hyperref[introduction-section-phantom]{Introduction}
\item \hyperref[conventions-section-phantom]{Conventions}
\item \hyperref[sets-section-phantom]{Théorie des ensembles}
\item \hyperref[categories-section-phantom]{Catégories}
\item \hyperref[topology-section-phantom]{Topologie}
\item \hyperref[sheaves-section-phantom]{Faisceaux sur les espaces}
\item \hyperref[sites-section-phantom]{Sites et faisceaux}
\item \hyperref[stacks-section-phantom]{Champs}
\item \hyperref[fields-section-phantom]{Corps}
\item \hyperref[algebra-section-phantom]{Algèbre commutative}
\item \hyperref[brauer-section-phantom]{Groupes de Brauer}
\item \hyperref[homology-section-phantom]{Algèbre homologique}
\item \hyperref[derived-section-phantom]{Catégories dérivées}
\item \hyperref[simplicial-section-phantom]{Méthodes simpliciales}
\item \hyperref[more-algebra-section-phantom]{Compléments d'algèbre}
\item \hyperref[smoothing-section-phantom]{Lissage des morphismes d'anneaux}
\item \hyperref[modules-section-phantom]{Faisceaux de modules}
\item \hyperref[sites-modules-section-phantom]{Modules sur les sites}
\item \hyperref[injectives-section-phantom]{Objets injectifs}
\item \hyperref[cohomology-section-phantom]{Cohomologie des faisceaux}
\item \hyperref[sites-cohomology-section-phantom]{Cohomologie sur les sites}
\item \hyperref[dga-section-phantom]{Algèbre différentielle graduée}
\item \hyperref[dpa-section-phantom]{Algèbre à puissances divisées}
\item \hyperref[sdga-section-phantom]{Faisceaux différentiels gradués}
\item \hyperref[hypercovering-section-phantom]{Hyperrecouvrements}
\end{enumerate}
Schémas
\begin{enumerate}
\setcounter{enumi}{25}
\item \hyperref[schemes-section-phantom]{Schémas}
\item \hyperref[constructions-section-phantom]{Constructions de schémas}
\item \hyperref[properties-section-phantom]{Propriétés des schémas}
\item \hyperref[morphisms-section-phantom]{Morphismes de schémas}
\item \hyperref[coherent-section-phantom]{Cohomologie des schémas}
\item \hyperref[divisors-section-phantom]{Diviseurs}
\item \hyperref[limits-section-phantom]{Limites de schémas}
\item \hyperref[varieties-section-phantom]{Variétés}
\item \hyperref[topologies-section-phantom]{Topologies sur les schémas}
\item \hyperref[descent-section-phantom]{Descente}
\item \hyperref[perfect-section-phantom]{Catégories dérivées de schémas}
\item \hyperref[more-morphisms-section-phantom]{Compléments sur les morphismes}
\item \hyperref[flat-section-phantom]{Compléments sur la platitude}
\item \hyperref[groupoids-section-phantom]{Schémas en groupoïdes}
\item \hyperref[more-groupoids-section-phantom]{Compléments sur les schémas en groupoïdes}
\item \hyperref[etale-section-phantom]{Morphismes étales de schémas}
\end{enumerate}
Thèmes de théorie des schémas
\begin{enumerate}
\setcounter{enumi}{41}
\item \hyperref[chow-section-phantom]{Homologie de Chow}
\item \hyperref[intersection-section-phantom]{Théorie de l'intersection}
\item \hyperref[pic-section-phantom]{Schémas de Picard des courbes}
\item \hyperref[weil-section-phantom]{Théories cohomologiques de Weil}
\item \hyperref[adequate-section-phantom]{Modules adéquats}
\item \hyperref[dualizing-section-phantom]{Complexes dualisants}
\item \hyperref[duality-section-phantom]{Dualité pour les schémas}
\item \hyperref[discriminant-section-phantom]{Discriminants et différentes}
\item \hyperref[derham-section-phantom]{Cohomologie de de Rham}
\item \hyperref[local-cohomology-section-phantom]{Cohomologie locale}
\item \hyperref[algebraization-section-phantom]{Géométrie algébrique et formelle}
\item \hyperref[curves-section-phantom]{Courbes algébriques}
\item \hyperref[resolve-section-phantom]{Résolution des surfaces}
\item \hyperref[models-section-phantom]{Réduction semi-stable}
\item \hyperref[functors-section-phantom]{Foncteurs et morphismes}
\item \hyperref[equiv-section-phantom]{Catégories dérivées de variétés}
\item \hyperref[pione-section-phantom]{Groupes fondamentaux de schémas}
\item \hyperref[etale-cohomology-section-phantom]{Cohomologie étale}
\item \hyperref[crystalline-section-phantom]{Cohomologie cristalline}
\item \hyperref[proetale-section-phantom]{Cohomologie pro-étale}
\item \hyperref[relative-cycles-section-phantom]{Cycles relatifs}
\item \hyperref[more-etale-section-phantom]{Compléments de cohomologie étale}
\item \hyperref[trace-section-phantom]{La formule des traces}
\end{enumerate}
Espaces algébriques
\begin{enumerate}
\setcounter{enumi}{64}
\item \hyperref[spaces-section-phantom]{Espaces algébriques}
\item \hyperref[spaces-properties-section-phantom]{Propriétés des espaces algébriques}
\item \hyperref[spaces-morphisms-section-phantom]{Morphismes d'espaces algébriques}
\item \hyperref[decent-spaces-section-phantom]{Espaces algébriques décents}
\item \hyperref[spaces-cohomology-section-phantom]{Cohomologie des espaces algébriques}
\item \hyperref[spaces-limits-section-phantom]{Limites d'espaces algébriques}
\item \hyperref[spaces-divisors-section-phantom]{Diviseurs sur les espaces algébriques}
\item \hyperref[spaces-over-fields-section-phantom]{Espaces algébriques sur des corps}
\item \hyperref[spaces-topologies-section-phantom]{Topologies sur les espaces algébriques}
\item \hyperref[spaces-descent-section-phantom]{Descente et espaces algébriques}
\item \hyperref[spaces-perfect-section-phantom]{Catégories dérivées d'espaces}
\item \hyperref[spaces-more-morphisms-section-phantom]{Compléments sur les morphismes d'espaces}
\item \hyperref[spaces-flat-section-phantom]{Platitude sur les espaces algébriques}
\item \hyperref[spaces-groupoids-section-phantom]{Groupoïdes dans les espaces algébriques}
\item \hyperref[spaces-more-groupoids-section-phantom]{Compléments sur les groupoïdes dans les espaces}
\item \hyperref[bootstrap-section-phantom]{Amorçage}
\item \hyperref[spaces-pushouts-section-phantom]{Sommes amalgamées d'espaces algébriques}
\end{enumerate}
Thèmes de géométrie
\begin{enumerate}
\setcounter{enumi}{81}
\item \hyperref[spaces-chow-section-phantom]{Groupes de Chow des espaces}
\item \hyperref[groupoids-quotients-section-phantom]{Quotients de groupoïdes}
\item \hyperref[spaces-more-cohomology-section-phantom]{Compléments sur la cohomologie des espaces}
\item \hyperref[spaces-simplicial-section-phantom]{Espaces simpliciaux}
\item \hyperref[spaces-duality-section-phantom]{Dualité pour les espaces}
\item \hyperref[formal-spaces-section-phantom]{Espaces algébriques formels}
\item \hyperref[restricted-section-phantom]{Algébrisation des espaces formels}
\item \hyperref[spaces-resolve-section-phantom]{Résolution des surfaces revisitée}
\end{enumerate}
Théorie des déformations
\begin{enumerate}
\setcounter{enumi}{89}
\item \hyperref[formal-defos-section-phantom]{Théorie formelle des déformations}
\item \hyperref[defos-section-phantom]{Théorie des déformations}
\item \hyperref[cotangent-section-phantom]{Le complexe cotangent}
\item \hyperref[examples-defos-section-phantom]{Problèmes de déformation}
\end{enumerate}
Champs algébriques
\begin{enumerate}
\setcounter{enumi}{93}
\item \hyperref[algebraic-section-phantom]{Champs algébriques}
\item \hyperref[examples-stacks-section-phantom]{Exemples de champs}
\item \hyperref[stacks-sheaves-section-phantom]{Faisceaux sur les champs algébriques}
\item \hyperref[criteria-section-phantom]{Critères de représentabilité}
\item \hyperref[artin-section-phantom]{Axiomes d'Artin}
\item \hyperref[quot-section-phantom]{Espaces de Quot et de Hilbert}
\item \hyperref[stacks-properties-section-phantom]{Propriétés des champs algébriques}
\item \hyperref[stacks-morphisms-section-phantom]{Morphismes de champs algébriques}
\item \hyperref[stacks-limits-section-phantom]{Limites de champs algébriques}
\item \hyperref[stacks-cohomology-section-phantom]{Cohomologie des champs algébriques}
\item \hyperref[stacks-perfect-section-phantom]{Catégories dérivées de champs}
\item \hyperref[stacks-introduction-section-phantom]{Introduction aux champs algébriques}
\item \hyperref[stacks-more-morphisms-section-phantom]{Compléments sur les morphismes de champs}
\item \hyperref[stacks-geometry-section-phantom]{La géométrie des champs}
\end{enumerate}
Thèmes de théorie des espaces de modules
\begin{enumerate}
\setcounter{enumi}{107}
\item \hyperref[moduli-section-phantom]{Champs de modules}
\item \hyperref[moduli-curves-section-phantom]{Espaces de modules de courbes}
\end{enumerate}
Divers
\begin{enumerate}
\setcounter{enumi}{109}
\item \hyperref[examples-section-phantom]{Exemples}
\item \hyperref[exercises-section-phantom]{Exercices}
\item \hyperref[guide-section-phantom]{Guide bibliographique}
\item \hyperref[desirables-section-phantom]{Desiderata}
\item \hyperref[coding-section-phantom]{Style de programmation}
\item \hyperref[obsolete-section-phantom]{Obsolète}
\item \hyperref[fdl-section-phantom]{Licence de documentation libre GNU}
\item \hyperref[index-section-phantom]{Index généré automatiquement}
\end{enumerate}
\end{multicols}


\bibliography{my}
\bibliographystyle{amsalpha}

\end{document}
